%! TeX root = ./main.tex
% vim: set spelllang=en:
\chapter{Phylogenetic reconciliation}\label{chap:reconciliation}

Two genes are said to be \emph{homologous} if they descend from a common ancestral gene.
When their divergence originates from a cospeciation with the ancestral host species, these genes are said to be \emph{orthologous}~\parencite{Fitch1970DistinguishingHomologousAnalogous}.
Conversely, the divergence can be independent of the host species: a duplication --- in which case the genes are referred to as \emph{paralogous} --- or a horizontal transfer --- they are then called \emph{xenologous}.
In particular, the identification of orthology relationships is an important step in the construction of species phylogenies, or in the study of gene function.
Indeed, orthologous genes more often share the same function than duplicated genes, which are subject to weaker selective pressure~\parencite{Lynch2001ProbabilityPreservationNewly}.

In the phylogeny of a family of homologous genes, the set of internal nodes can be partitioned according to the type of evolutionary event to which they correspond.
This partitioning makes it possible to qualify any pair of genes in the tree as orthologous, paralogous or xenologous.
The so-called \emph{reconciliation} method, introduced by \textcite{Goodman1979FittingGeneLineage}, formalized by \textcite{Page1994MapsTreesCladistic} and subsequently generalized in several ways, makes it possible to carry out such a partitioning.
It is based on the fact that genes and their host species form an evolutionary hierarchy, since the evolution of hosts affects that of the genes they contain.

In this chapter, we first introduce the reconciliation of a single family of units, which we call \emph{bilateral}, and then the various generalizations that allow \emph{multiple} reconciliations of several interdependent units.

%! TeX root = ./main.tex
% vim: set spelllang=en:
\section{Bilateral reconciliation}\label{sec:reconciliation:bilateral}

\subsection{Embedding phylogenetic trees}

An internal node~\(v\) in a phylogenetic tree~\(T\) of evolutionary units can be caused either by a codivergence with the associated host, or by a divergence independent of the host.
When~\(v\) results from a codivergence, it corresponds to an internal node in the phylogeny~\(S\) of the hosts.
Conversely, if~\(v\) is independent of the evolution of the hosts, it corresponds to a point along one of the branches of \(S\).
Finally, the leaves of \(T\) represent extant units and each correspond to a leaf of~\(S\), their host.
Such an association of \(T\) with \(S\) is called a \emph{reconciliation}.

Formally, a reconciliation is an embedding of \(T\) into \(S\), in the sense defined below.
To simplify the definition, when \(v\) corresponds to a branch of \(S\), it is associated with its endpoint.

\begin{definition}[Reconciliations]\label{definition:rec}
    Let \(T\) and \(S\) be two phylogenetic trees and \(\llmap : \tleavs(T) \onto \tleavs(S)\) a map that associates each leaf of \(T\) with its host in \(S\).
    A \emph{reconciliation} of \(T\) with~\(S\) is a map \(\rec : \tnodes(T) \onto \tnodes(S)\) that extends \(\llmap\) such that no pair of nodes \((v, w) \in \tnodes(T)^2\) is such that \(v \tancsneq w\) and \(\rec(v) \tdescneq \rec(w)\).
\end{definition}

Two examples of maps are illustrated in \cref{fig:reconciliation:strict-rec}.
The map \(\rec_1\) is an acceptable reconciliation, whereas the map \(\rec_2\) does not satisfy the condition of the definition that forbids reconciliations from \enquote{going back in time} with respect to the topology of \(S\).

\begin{figure}[htb]
    \capstart
    \centering
    %! TeX root = ../main.tex
% vim: set spelllang=fr:
\begin{tikzpicture}
    %! TeX root = ../main.tex
% vim: set spelllang=fr:

\tikzset{
    % Predefined dependent tree
    graphs/declare={ttree}{
        l3s1 / [coordinate] -- {
            l2s2 / [coordinate] -- {
                l1s4 / 3,
                l1s3 / 1,
            },
            l2s1 / [coordinate] -- {
                l1s2 / 2,
                l1s1 / 1,
            },
        },
    },
    % Predefined host tree
    graphs/declare={stree}{
        l3s1 / [coordinate] -- {
            l1s3 / 3 [> length=2],
            l2s1 / [coordinate] -- {
                l1s2 / 2,
                l1s1 / 1,
            },
        },
    },
}


    \graph[phylogenetic tree layout, nodes={font=\small\itshape}]{
        ttree[name=ct]
    };
    \node at (0, -2.3) {\(T\)};

    \begin{scope}[shift={(3.25,0)}]
        \graph[phylogenetic tree layout, nodes={font=\small\bfseries}]{
            stree[name=ls]
        };
        \coordinate (labels) at (0, -2.3);
        \node at (labels-|ls l1s2) {\(S\)};
    \end{scope}

    \draw[separator rule] (5.15, .15) -- ++(0, -2.6);

    \begin{scope}[shift={(7.5,0)}]
        \graph[phylogenetic tree layout, nodes={font=\small\itshape}]{
            ttree[name=lt]
        };
        \node at (0, -2.3) {\(T\)};

        \begin{scope}[shift={(3.25,0)}]
            \graph[phylogenetic tree layout, nodes={font=\small\bfseries}]{
                stree[name=rs]
            };
            \coordinate (labels) at (0, -2.3);
            \node at (labels-|rs l1s2) {\(S\)};
        \end{scope}
    \end{scope}

    \path[overlaid associate=reccolor1] (ct l2s1) edge[bend left=15] (ls l2s1);
    \path[overlaid associate=reccolor1] (ct l2s2) edge[out=0, in=180] (ls l3s1);
    \path[associate=reccolor1] (ct l3s1) edge[bend left=10]
        node[above, midway] {\(\rec_1\)} (ls l3s1);

    \path[overlaid associate=reccolor2] (lt l3s1) edge[bend left=10] (rs l2s1);
    \path[overlaid associate=reccolor2] (lt l2s1) edge[bend left=15] (rs l2s1);
    \path[overlaid associate=reccolor2] (lt l2s2) edge[out=0, in=180]
        node[above, midway, reccolor2] {\(\rec_2\)} (rs l3s1);
\end{tikzpicture}

    \caption[Two examples of strict reconciliations]{
        Two associations of the phylogeny \(T\) with \(S\).
        The associations between internal nodes are represented by arrows and the associations between leaves are indicated implicitly by assigning them the same label.
    }
    \label{fig:reconciliation:strict-rec}
\end{figure}

Some units are able to migrate between two coexisting hosts.
For example, particularly in prokaryotes, horizontal transfers can cause genes to migrate between species.
In a reconciliation \(\rec\), a branch \((v, w) \in \tedges(T)\) such that \(\rec(v) \nottcomp \rec(w)\) implies that a host switch took place between two points on the path between \(\rec(v)\) and \(\rec(w)\).
Note, however, that \emph{unobservable} host switches can take place.
For example if, between \(v\) and \(w\), the unit migrated from host \(\rec(v)\) to a separated host \(\gamma\), then back again to a host descending from \(\rec(v)\) without leaving any trace among the descendants of \(\gamma\), the branch \((v, w)\) will not appear as a host-switch branch.

We call a reconciliation that does not contain any such branch \emph{strict}, since the evolution of the units then strictly follows that of their hosts.
In \cref{fig:reconciliation:strict-rec}, the reconciliation \(\rec_1\) is a strict reconciliation.
\Cref{fig:reconciliation:free-rec} gives two examples of free reconciliations.

\begin{definition}\label{definition:rec-strict}
    A reconciliation \(\rec\) of \(T\) with \(S\) is said to be \emph{strict} if all pairs of nodes
    \((v, w) \in \tnodes(T)^2\) with \(v \tancs w\) are such that \(\rec(v) \tancs \rec(w)\).
    Otherwise, \(\rec\) is said to be \emph{free}.
\end{definition}

\begin{figure}[htb]
    \capstart
    \centering
    %! TeX root = ../main.tex
% vim: set spelllang=fr:
\begin{tikzpicture}
    %! TeX root = ../main.tex
% vim: set spelllang=fr:

\tikzset{
    % Predefined dependent tree
    graphs/declare={ttree}{
        l3s1 / [coordinate] -- {
            l2s2 / [coordinate] -- {
                l1s4 / 3,
                l1s3 / 1,
            },
            l2s1 / [coordinate] -- {
                l1s2 / 2,
                l1s1 / 1,
            },
        },
    },
    % Predefined host tree
    graphs/declare={stree}{
        l3s1 / [coordinate] -- {
            l1s3 / 3 [> length=2],
            l2s1 / [coordinate] -- {
                l1s2 / 2,
                l1s1 / 1,
            },
        },
    },
}


    \graph[phylogenetic tree layout, simple, nodes={font=\small\itshape}]{
        ttree[name=t1];
        t1 l2s2 --[transfer branch] t1 l1s4;
    };
    \node at (0, -2.3) {\(T\)};

    \begin{scope}[shift={(3.25, 0)}]
        \graph[phylogenetic tree layout, nodes={font=\small\bfseries}]{
            stree[name=s]
        };
        \coordinate (labels) at (0, -2.3);
        \node at (labels-|s l1s2) {\(S\)};
    \end{scope}

    \path[overlaid associate=reccolor3] (t1 l2s1) edge[bend left=13] (s l2s1);
    \path[associate=reccolor3] (t1 l3s1) edge[bend left=10]
        node[above, midway] {\(\rec_3\)} (s l2s1);
    \path[associate=reccolor3] (t1 l2s2) edge[bend left=10, shorten >=0pt] (s l1s1);

    \draw[separator rule] (4.8, .15) -- ++(0, -2.6);

    \begin{scope}[shift={(7.5, 0)}]
        \graph[
            phylogenetic tree layout,
            evolutionary unit length=0.44cm,
            nodes={font=\small\itshape}
        ]{
            t l4s1 / [
                coordinate,
                label={[host, reccolor4, yshift=3pt]left:α},
            ] -- {
                t l3s2 / [
                    coordinate,
                    label={[host, reccolor4, yshift=3pt]right:1},
                ] -- {
                    t l2s2 / [
                        coordinate,
                        label={[host, reccolor4]right:γ},
                        > transfer branch,
                    ] -- {
                        t l1s6 / 4 -- {},
                        t l1s5 / 3 -- {},
                    },
                    t l1s4 / 1 [> length=2],
                },
                t l3s1 / [
                    coordinate,
                    label={[host, reccolor4, yshift=3pt]left:4}
                ] -- {
                    t l1s3 / 4 [> length=2],
                    t l2s1 / [
                        coordinate,
                        label={[host, reccolor4]left:β},
                        > transfer branch,
                    ] -- {
                        t l1s2 / 2 -- {},
                        t l1s1 / 1 -- {},
                    },
                },
            },
        };
        \node[reccolor4] at (2, 0.15) {\((\rec_4)\)};
        \node at (0, -2.3) {\(T'\)};

        \begin{scope}[shift={(3.75, 0)}]
            \graph[phylogenetic tree layout, nodes={font=\small\bfseries}]{
                s l3s1 / [coordinate, label={[host]left:α}] -- {
                    s l2s2 / [coordinate, label={[host]right:γ}] -- {
                        s l1s4 / 4,
                        s l1s3 / 3,
                    },
                    s l2s1 / [coordinate, label={[host]left:β}] -- {
                        s l1s2 / 2,
                        s l1s1 / 1,
                    },
                },
            };
            \node at (0, -2.3) {\(S'\)};
        \end{scope}
    \end{scope}
\end{tikzpicture}

    \caption[Two examples of free reconciliations]{
        Two free reconciliations of the phylogenies \(T\) with \(S\) and \(T'\) with \(S'\).
        For the sake of readability, \(\rec_4\) is represented by label equalities rather than by arrows.
        The branches marked with dotted lines correspond to host switches.
    }
    \label{fig:reconciliation:free-rec}
\end{figure}

Some authors restrict \cref{definition:rec} to forbid a given internal node from having its two children associated with two separated hosts~\parencite{Nojgaard2018TimeConsistentReconciliation}.
Adding this condition guarantees that there exists a path without host switches from any internal node to a leaf.
Thus, every internal node corresponds to a divergence of which there is an extant trace.
The models we consider hereafter, however, require keeping the generality of the definition.

\subsection{Reconciled trees and reconciliation models}\label{sec:reconciliation:bilateral:model}

A reconciliation model specifies the types of events that can affect, jointly or separately, the evolution of the units and of their hosts.
An interpretation of a reconciliation under a given model is a combination of events from the model that explains its unit--host associations.
These interpretations are represented by reconciled trees labelled with the chosen events (see examples in \cref{fig:reconciliation:rec-trees}).

\begin{definition}[Reconciled trees]\label{def:reconciliation:rectree}
    Let \(T\) and \(S\) be two trees, \(\rec\) a reconciliation of \(T\) with \(S\), and \(\recmodel\) a set of event types.
    A \emph{reconciled tree} \(\rectree_{\!\recmodel}(\rec)\) (or simply \(\rectree(\rec)\) when the model used is clear from the context) is a partially binary extension of~\(T\) equipped with a map \(\exrec : \tnodes(\rectree) \onto \tnodes(S)\) that extends \(\rec\) and with a map \(\evt : \tnodes(\rectree) \onto \recmodel \cup \{\evtE\}\)
    such that \(\tleavs(T) = \{ v \in \tleavs(\rectree) \suchthat \evt(v) = \evtE \}\).
\end{definition}

Two main classes of methods exist to assess the plausibility of a reconciliation under a given model.
In \emph{parsimony} methods, the reconciliations considered most plausible are those whose interpretation invokes the fewest events of the model.
In the \emph{probabilistic} framework, it is the reconciliations whose best interpretation maximizes the likelihood that are favoured~\parencite{Doyon2011ModelsAlgorithmsPrograms}.
This thesis focuses on parsimony methods.

In parsimony, a function \(\cost : \recmodel \onto \realpos\) assigning a cost to each event type is fixed in advance.
The cost of a reconciled tree is the sum of the costs of its events.
The cost of a reconciliation is then defined as that of its most economical interpretation.
This definition gives rise to a family of combinatorial optimization problems that consist in searching for the reconciliations of minimum cost, that is, the most parsimonious ones.

\begin{problem}{\(\recmodel\)-parsimony\label{prob:reconciliation:parsimony}}
    \probinput{Phylogenies \(T\) and \(S\);
               leaf association \(\llmap : \tleavs(T) \onto \tleavs(S)\);
               costs \(\cost : \recmodel \onto \realpos\).}
    \proboutput{Minimum cost of a reconciliation of \(T\) with \(S\) under the model \(\recmodel\) and the costs \(\cost\).}
\end{problem}

Even if the phylogenies \(T\) and \(S\) under study are free of errors, the most parsimonious reconciliations do not necessarily correspond to complete or exact histories of the units concerned, since evolution rarely takes the shortest path.
On the other hand, this method has the advantage of selecting reconciliations that explain the observed data while requiring the fewest additional assumptions.

\begin{figure}[htb]
    \capstart
    \centering
    %! TeX root = ../main.tex
% vim: set spelllang=fr:
\begin{tikzpicture}
    %% Réconciliation DL
    \begin{scope}[x=.5cm, y=.8cm]
        \graph[no placement, nodes={dependent}]{
            t1l4s1 / [duplication, y=-.16] -- {
                t1l3s1 / [speciation, x=-.5, y=-1, outer sep=-.5pt] -- {
                    t1l2s1 / [speciation, x=-2.5, y=-2, outer sep=-.8pt] -- {
                        t1l1s1 / 1 [x=-4, y=-3],
                        t1l1s2 / 2 [x=-1, y=-3],
                    },
                    t1l1s3 / [loss, >loss branch, rotate=35, x=2, y=-3],
                },
                t1l3s2 / [speciation, x=.5, y=-1, outer sep=-.5pt] -- {
                    t1l2s2 / [speciation, x=-1.5, y=-2, outer sep=-.8pt] -- {
                        t1l1s4 / 1 [x=-3, y=-3],
                        t1l1s5 / [loss, >loss branch, rotate=40, x=0, y=-3],
                    },
                    t1l1s6 / 3 [x=3, y=-3, outer sep=0pt],
                },
            },
        };
        \node at (-6, -1.5) {\(\rectreeDL(\rec_1)\)};
    \end{scope}

    \draw[separator rule] (2.75, 0) -- ++(0, -2.4);

    %% Réconciliation DTL
    \begin{scope}[x=.5cm, y=.8cm, shift={(12, 0)}]
        \graph[no placement, nodes={dependent}]{
            t2l4s1 / [duplication, x=-1, y=-.75] -- {
                t2l3s1 / [speciation, x=-2.5, y=-1.5, outer sep=-.7pt] -- {
                    t2l1s1 / 1 [x=-4, y=-3],
                    t2l1s2 / 2 [x=-1, y=-3],
                },
                t2l3s2 / [speciation, x=-1.5, y=-1.5, outer sep=-.4pt] -- {
                    t2l2s1 / [transfer, x=-2.4, y=-2.4] -- {
                        t2l1s3 / 1 [x=-3, y=-3],
                        t2l1s4 anchor / [coordinate, x=1.95, y=-2.4, > transfer branch]
                            -- t2l1s4 / 3 [x=2.5, y=-3],
                    },
                    t2l1s5 / [loss, >loss branch, rotate=32, x=0, y=-3],
                },
            },
        };
        \coordinate (t2lroot) at (0, -.16);
        \node at (6, -1.5) {\(\rectreeDTL(\rec_3)\)};
    \end{scope}

    \begin{pgfonlayer}{hostnodes}
        \node[host leaf node, fit=(t1l1s1)(t1l1s4)] (s1l1s1) {};
        \node[host leaf node, fit=(t1l1s2)(t1l1s5)] (s1l1s2) {};
        \node[host leaf node, fit=(t1l1s3)(t1l1s6)] (s1l1s3) {};
        \node[host node, fit=(t1l2s1)(t1l2s2)] (s1l2s1) {};
        \node[host node, fit=(t1l3s1)(t1l3s2)] (s1l3s1) {};

        \node[host leaf node, fit=(t2l1s1)(t2l1s3)] (s2l1s1) {};
        \node[host leaf node, fit=(t2l1s2)(t2l1s5)] (s2l1s2) {};
        \node[host leaf node, fit=(t2l1s4), inner xsep=.2cm] (s2l1s3) {};
        \node[host node, fit=(t2l3s1)(t2l3s2)] (s2l2s1) {};
        \node[host node, fit=(t2lroot), inner xsep=.5cm, inner ysep=.25cm] (s2l3s1) {};
    \end{pgfonlayer}

    \foreach \i in {1,2} {
        \begin{pgfonlayer}{hostedges}
            \path[host edge] (s\i l2s1.east) -- (s\i l3s1.east) -- (s\i l3s1.west) -- (s\i l2s1.160);
            \path[host edge] (s\i l1s1.-10) -- (s\i l2s1.east) -- (s\i l2s1.160) -- (s\i l1s1.165);
            \path[host edge] (s\i l1s2.15) -- (s\i l2s1.east) -- (s\i l2s1.west) -- (s\i l1s2.190);
            \path[host edge] (s\i l1s3.15) -- (s\i l3s1.east) -- (s\i l3s1.west) -- (s\i l1s3.190);
        \end{pgfonlayer}
    };

    \begin{pgfonlayer}{hostedges}
        \path[host edge terminal] (s1l3s1.east) -- ($(s1l3s1.east)+(0, 1)$)
            -- ($(s1l3s1.west)+(0, 1)$) -- (s1l3s1.west);
    \end{pgfonlayer}
\end{tikzpicture}

    \caption[Two examples of reconciled trees]{
        Reconciled trees for the reconciliations \(\rec_1\) and \(\rec_3\) (from \cref{fig:reconciliation:strict-rec,fig:reconciliation:free-rec}) under the duplication-loss~(\(\recmodelDL\)) and duplication-transfer-loss~(\(\recmodelDTL\)) models respectively.
        Each tree is overlaid on the outline of the host phylogeny so as to reflect the association of nodes given by the reconciliation.
        Each circle~{\showspeciation} represents a codivergence event~(\(\evtS\)),
        each square~{\showduplication} a duplication~(\(\evtD\)),
        each diamond~{\showtransfer} a transfer~(\(\evtT\)),
        and each upward-opening half-circle~{\showloss} a loss~(\(\evtL\)).
    }
    \label{fig:reconciliation:rec-trees}
\end{figure}

\paragraph{Duplication and loss~\textup{\parencite{Goodman1979FittingGeneLineage,Page1994MapsTreesCladistic}}}
This model comprises three event types, \(\recmodelDL = \{\evtD, \evtL, \evtS\}\).
The divergence of a unit is explained by its duplication~(\(\evtD\)) or its cospeciation with its host~(\(\evtS\)), and the absence of a unit in a descendant host is explained by a loss~(\(\evtL\)).
Each node \(v \in \tnodes(\rectree)\) of a reconciled tree under this model is such that:
\begin{compactitemize}
\item if \(\evt(v) = \evtD\), then \(v\) is an internal node and \(\exrec(v) = \exrec(v\lchild) = \exrec(v\rchild)\);
    \item if \(\evt(v) = \evtL\), then \(v\) is a leaf;
    \item if \(\evt(v) = \evtS\), then \(v\) and \(\exrec(v)\) are internal nodes and \(\{\exrec(v\lchild), \exrec(v\rchild)\} = \{\exrec(v)\lchild, \exrec(v)\rchild\}\).
\end{compactitemize}
Loss events introduce additional leaves into the reconciled tree compared to the tree of units.
They make it possible to interpret situations where the two endpoints of a branch are associated with two non-adjacent hosts, by inserting a sequence of speciations and losses into this branch.
\Cref{fig:reconciliation:rec-trees} presents an example of an interpretation of a reconciliation under the \(\recmodelDL\) model in which such a sequence of losses is used.

The cost functions considered for this model usually set \(\cost(\evtS) = 0\).
Consequently, searching for minimum-cost reconciliations amounts to maximizing the number of codivergences between \(T\) and \(S\).
This choice is motivated by the fact that each \(\evtS\)~event relies on the speciation of a host, an event already postulated by the tree \(S\) given as input.

\paragraph{Host switch~\textup{\parencite{Tofigh2011SimultaneousIdentificationDuplications,Ranwez2016InferringGeneDuplications}}}
The model \(\recmodelDTL = \{\evtD, \evtT, \evtL, \evtS\}\) extends the \(\recmodelDL\) model to integrate the possibility of host switches and thus allow free reconciliations to be interpreted.
In addition to the three event types defined previously, for an internal node \(v \in \tinner(\rectree)\) under this model, if \(\evt(v) = \evtT\), then \(\exrec(v) = \exrec(v\lchild)\) and \(\exrec(v) \nottcomp \exrec(v\rchild)\), or else \(\exrec(v) = \exrec(v\rchild)\) and \(\exrec(v) \nottcomp \exrec(v\lchild)\).
An example of an interpretation of a reconciliation under the \(\recmodelDTL\) model is given in \cref{fig:reconciliation:rec-trees}.

In this model, a host switch can be immediately followed by a loss, thus leaving no trace in the original host.
This observation, initially ignored in the work of \textcite{Tofigh2011SimultaneousIdentificationDuplications}, can lead to the construction of significantly different interpretations \parencite{Ranwez2016InferringGeneDuplications}.
It also makes it possible to interpret reconciliations containing nodes whose two children are associated with hosts separated from that of their parent.
However, taking it into account leads to an increase in the number of co-optimal solutions in the associated optimization problem~\parencite{Weiner2021ImprovedDuplicationTransfer}.

\paragraph{Unsampled hosts~\textup{\parencite{Szoelloesi2013LateralGeneTransfer,Weiner2021ImprovedDuplicationTransfer}}}
The phylogeny \(S\) is necessarily incomplete since it only gives information about the hosts of the units of~\(T\).
The internal nodes leading only to hosts absent from this sample are therefore missing.
In the context of free reconciliations, it is possible for a unit to evolve in a part of the phylogeny located outside the sample before later being transferred to a sampled host.
To represent this possibility, the \(\recmodelDTLx\) model adds an additional event, denoted \(\evtTX\), to the \(\recmodelDTL\) model.
For an internal node \(v \in \tinner(\rectree)\) under this model, if \(\evt(v) = \evtTX\), then \(\exrec(v) = \exrec(v\lchild)\) and \(\exrec(v) \tancs \exrec(v\rchild)\), or else \(\exrec(v) = \exrec(v\rchild)\) and \(\exrec(v) \tancs \exrec(v\lchild)\).
This event offers an alternative to inserting sequences of losses to interpret a reconciliation branch whose two endpoints are associated with non-adjacent hosts.
This can provide a more parsimonious interpretation for some reconciliations, as illustrated in~\cref{fig:reconciliation:rec-dtlx}.

\begin{figure}[htb]
    \capstart
    \centering
    %! TeX root = ../main.tex
% vim: set spelllang=fr:
\begin{tikzpicture}
    %% Réconciliation DL
    \begin{scope}[x=.5cm, y=.8cm]
        \graph[no placement, nodes={dependent}]{
            t1l4s1 / [duplication, y=-.16] -- {
                t1l3s1 / [speciation, x=-.5, y=-1, outer sep=-.45pt] -- {
                    t1l2s1 / [speciation, x=-2.5, y=-2, outer sep=-.8pt] -- {
                        t1l1s1 / 1 [x=-4, y=-3, >shorten >=-3pt],
                        t1l1s2 / 2 [x=-1, y=-3, >shorten >=-4pt],
                    },
                    t1l1s3 / 3 [x=2, y=-3, >shorten >=-3pt],
                },
                t1l3s2 / [speciation, x=.5, y=-1, outer sep=-.45pt] -- {
                    t1l2s2 / [speciation, x=-1.5, y=-2, outer sep=-.8pt] -- {
                        t1l1s4 / 1 [x=-3, y=-3, >shorten >=-3pt],
                        t1l1s5 / [loss, >loss branch, rotate=40, x=0, y=-3],
                    },
                    t1l1s6 / [loss, >loss branch, rotate=40, x=3, y=-3],
                },
            },
        };
    \end{scope}

    \draw[separator rule] (2.75, 0) -- ++(0, -2.4);

    %% Réconciliation DTL
    \begin{scope}[x=.5cm, y=.8cm, shift={(12, 0)}]
        \graph[no placement, nodes={dependent}]{
            t2l4s1 / [speciation, x=0.5, y=-.16, outer sep=-.7pt] -- {
                t2l3s1 / [speciation, x=-.8, y=-.9, outer sep=-.7pt] -- {
                    t2l2s1 / [speciation, x=-2.25, y=-1.75, outer sep=-.7pt] -- {
                        t2l1s1 / 1 [x=-4, y=-3, >shorten >=-3pt],
                        t2l1s2 / 2 [x=-.5, y=-3, >shorten >=-4pt],
                    },
                    t2l1s3 / 3 [x=2, y=-3, >shorten >=-3pt],
                },
                t2l3s2 / [transfer, x=3.75, y=-2.4] -- {
                    t2l2s2 / [coordinate, x=-2.4, y=-2.4, > transfer branch] --
                        t2l1s4 / 1 [x=-3, y=-3, >shorten >=-3pt],
                    t2l1sx / [coordinate, x=5, y=-3.3, > dotted],
                },
            },
        };
    \end{scope}

    \begin{pgfonlayer}{hostnodes}
        \node[host leaf node, fit=(t1l1s1)(t1l1s4)] (s1l1s1) {};
        \node[host leaf node, fit=(t1l1s2)(t1l1s5)] (s1l1s2) {};
        \node[host leaf node, fit=(t1l1s3)(t1l1s6)] (s1l1s3) {};
        \node[host node, fit=(t1l2s1)(t1l2s2)] (s1l2s1) {};
        \node[host node, fit=(t1l3s1)(t1l3s2)] (s1l3s1) {};

        \node[host leaf node, fit=(t2l1s1)(t2l1s4)] (s2l1s1) {};
        \node[host leaf node, fit=(t2l1s2), inner xsep=.075cm] (s2l1s2) {};
        \node[host leaf node, fit=(t2l1s3), inner xsep=.075cm] (s2l1s3) {};
        \node[host node, fit=(t2l2s1), inner xsep=.15cm] (s2l2s1) {};
        \node[host node, fit=(t2l3s1), inner xsep=.15cm] (s2l3s1) {};
        \node[host node, fit=(t2l4s1), inner xsep=.15cm] (s2l4s1) {};
        \node[ellipse, fit=(t2l1sx), inner sep=.3cm] (s2l1sx) {};
    \end{pgfonlayer}

    \begin{pgfonlayer}{hostedges}
        \path[host edge] (s1l2s1.east) -- (s1l3s1.east) -- (s1l3s1.west) -- (s1l2s1.160);
        \path[host edge] (s1l1s1.-10) -- (s1l2s1.east) -- (s1l2s1.160) -- (s1l1s1.165);
        \path[host edge] (s1l1s2.15) -- (s1l2s1.east) -- (s1l2s1.west) -- (s1l1s2.190);
        \path[host edge] (s1l1s3.15) -- (s1l3s1.east) -- (s1l3s1.west) -- (s1l1s3.190);
        \path[host edge terminal] (s1l3s1.east) -- ($(s1l3s1.east)+(0, 1)$)
            -- ($(s1l3s1.west)+(0, 1)$) -- (s1l3s1.west);

        \path[host edge unsampled] (s2l4s1.15) -- (s2l1sx.east) -- (s2l1sx.west) -- (s2l4s1.west);
        \path[host edge] (s2l3s1.-10) -- (s2l4s1.-10) -- (s2l4s1.160) -- (s2l3s1.160);
        \path[host edge] (s2l2s1.-10) -- (s2l3s1.-10) -- (s2l3s1.160) -- (s2l2s1.160);
        \path[host edge] (s2l1s1.-10) -- (s2l2s1.east) -- (s2l2s1.160) -- (s2l1s1.165);
        \path[host edge] (s2l1s2.15) -- (s2l2s1.east) -- (s2l2s1.west) -- (s2l1s2.190);
        \path[host edge] (s2l1s3.15) -- (s2l3s1.east) -- (s2l3s1.west) -- (s2l1s3.190);
    \end{pgfonlayer}
\end{tikzpicture}

    \caption[Alternative interpretations of the same reconciliation under two models]{
        Two alternative interpretations of the same reconciliation under the \(\recmodelDL\) and \(\recmodelDTLx\) models.
        In the reconciled tree on the right, the event at the root is a \(\evtTX\) event, which is represented as a speciation followed by a transfer inside a hypothetical unsampled (hatched) branch of the host phylogeny.
    }
    \label{fig:reconciliation:rec-dtlx}
\end{figure}

\paragraph{Other models}
Extensions of reconciliation to take \emph{incomplete lineage sorting}~(ILS) into account have been proposed by \textcite{Rasmussen2012UnifiedModelingGene,Wu2013MostParsimoniousReconciliation,Chan2017InferringIncompleteLineage}.
This phenomenon occurs when speciations that are close in time cause the emerging species to inherit distinct subsets of units.
As an alternative to \(\recmodelDL\)-type models, \textcite{Goodman1979FittingGeneLineage} proposed a model integrating the possibility for a gene to be activated or deactivated by a regulatory mechanism, which makes it possible to limit the inference of superfluous losses.
Finally, the dating of the host phylogeny, which imposes an additional constraint on host switches, can be integrated into the models allowing this type of event~\parencite{Tofigh2011SimultaneousIdentificationDuplications}.

\subsection{Time consistency}

Host switches imply temporal constraints, since the donor host and the recipient host of each switch must have been able to coexist for the event to occur.
Thus, for a reconciled tree to represent a realistic evolutionary scenario, it must satisfy the following constraints, imposed by the topology of the host tree and by the events of the reconciliation model
\parencite{Stolzer2012InferringDuplicationsLosses,Ma2017DTLReconciliationRepair}.

\begin{definition}[Time consistency of a reconciled tree]\label{def:reconciliation:time-cons}
    A reconciled tree \(\rectree(\rec)\) is said to be \emph{time-consistent} if there exists a preorder \(\succeq_\tau\) on \(\tnodes(\rectree) \cup \tnodes(S)\) such that:
    \begin{condition}
        \item \(s \succ_\tau s\child\) for each branch \((s, s\child) \in \tedges(S)\); \label{def:reconciliation:time-cons:hosts}
        \item \(v \succeq_\tau v\child\) for each branch \((v, v\child) \in \tedges(\rectree)\); \label{def:reconciliation:time-cons:events}
        \item \(\tparent(\exrec(v)) \succ_\tau v \succeq_\tau \exrec(v)\) for each node \(v \in \tnodes(\rectree)\) such that \(\exrec(v) \neq \troot(S)\), otherwise \(v \succeq_\tau \exrec(v)\); \label{def:reconciliation:time-cons:events-hosts}
        \item \(v \succeq_\tau v\child\) and \(v\child \succeq_\tau v\) for each branch \((v, v\child) \in \tedges(\rectree)\) such that \(\exrec(v) \nottcomp \exrec(v\child)\). \label{def:reconciliation:time-cons:host-switch} \qedhere
    \end{condition}
\end{definition}

Intuitively, a reconciled tree is time-consistent if it is possible to draw a diagram of it similar to those of \cref{fig:reconciliation:rec-trees,fig:reconciliation:rec-dtlx}.
In such a representation, all branches are oriented along the same time axis (\cref{def:reconciliation:time-cons:hosts,def:reconciliation:time-cons:events}).
Each codivergence is overlaid on the node of the corresponding host tree and each divergence is placed along the corresponding branch (\cref{def:reconciliation:time-cons:events-hosts}).
Finally, host-switch branches are drawn orthogonally to the time axis (\cref{def:reconciliation:time-cons:host-switch}).

\begin{figure}[htb]
    \capstart
    \centering
    \begin{tikzpicture}[scale=1.25]
    \begin{scope}
        \tikzset{
    x={(1pt,0)},
    y={(0,-1pt)},
    limit width/.style 2 args={
        execute at begin node={\begin{varwidth}{#1}#2},
        execute at end node={\end{varwidth}},
    },
    host background/.style={
        fill=host background color,
        draw=host background color,
        line width={.01pt},
    },
    unsampled host background/.style={
        preaction={fill=host background color},
        draw=host background color,
        pattern={Lines[angle=45, distance=3pt, line width=1.5pt]},
        pattern color=white,
    },
    host label/.style={
        font=\bfseries,
        midway,
        anchor=north,
        limit width={12em}{\centering},
    },
    branch/.style={
        draw={#1},
        line width={0.5pt},
    },
    loss branch/.style={
        branch,
        draw={black!40},
    },
    transfer branch/.style={
        branch={#1},
        densely dashed,
        -{Latex[length=0pt 8]},
    },
    invalid transfer branch/.style={
        transfer branch={red!80!black},
        out distance=20,
        in distance=20,
    },
    event/.style={
        draw={#1}, fill={host background color!50!white},
        font={\color{#1}},
        outer sep=0pt, inner xsep=0pt, inner ysep=2pt,
        limit width={6em}{\narrowragged},
        line width={0.5pt},
    },
    event/.default={black},
    extant/.code 2 args={
        \pgfkeysalso{
            circle, fill={#1},
            outer sep=0pt, inner sep=0pt,
            minimum size={3},
        }
        \def\cmp{#2}\ifx\cmp\empty\else
            \pgfkeysalso{
                label={[%
                    font={\color{#1}\vphantom{gb}},
                    fill=host background color!35!white,
                    rounded corners=2,
                    label distance=2pt,
                    inner xsep=2pt, inner ysep=1pt,
                    outer xsep=0pt, outer ysep=0pt,
                    limit width={6em}{\centering},
                ]below:#2},
            }
        \fi
    },
    extant/.default={black}{},
    unsampled/.style={
        circle, fill={#1},
        outer sep=0pt, inner sep=0pt,
        minimum size={3},
    },
    unsampled/.default={black!40},
    speciation/.style 2 args={
        event={#1}, rectangle, rounded corners,
        inner xsep=4pt,
        minimum width={8},
        minimum height={8},
        execute at begin node={#2},
    },
    duplication/.style 2 args={
        event={#1}, rectangle,
        inner xsep=4pt,
        minimum width={8},
        minimum height={8},
        execute at begin node={#2},
    },
    transfer/.code 2 args={
        \pgfkeysalso{
            event={#1},
        }
        \def\cmp{#2}\ifx\cmp\empty
            \pgfkeysalso{
                diamond,
                inner sep=1.5pt,
            }
        \else
            \pgfkeysalso{
                chamfered rectangle,
                chamfered rectangle sep={8 / 2.4},
                inner xsep=2pt,
                inner ysep=-1pt,
                minimum width={8},
                minimum height={8},
                execute at begin node={#2},
            }
        \fi
    },
    gain/.code 2 args={
        \pgfkeysalso{
            draw={#1}, fill=black,
            semicircle, inner sep=1pt,
            shape border rotate={0},
        }
        \def\cmp{#2}\ifx\cmp\empty\else
            \pgfkeysalso{
                label={[%
                    font={\scriptsize\vphantom{gb}},
                    label distance=2pt,
                    inner xsep=0pt, inner ysep=0pt,
                    outer xsep=0pt, outer ysep=0pt,
                    limit width={6em}{\narrowragged},
                ]right:#2},
            }
        \fi
    },
    loss/.code 2 args={
        \pgfkeysalso{
            draw={#1}, fill=black,
            semicircle, inner sep=1pt,
            shape border rotate={180+0},
        }
        \def\cmp{#2}\ifx\cmp\empty
            \pgfkeysalso{draw={black!40}, fill={black!40}}
        \else
            \pgfkeysalso{
                label={[%
                    font={\scriptsize\vphantom{gb}},
                    label distance=2pt,
                    inner xsep=0pt, inner ysep=0pt,
                    outer xsep=0pt, outer ysep=0pt,
                    limit width={6em}{\narrowragged},
                ]right:#2},
            }
        \fi
    },
}
% hosts
\path[host background] (0.0,-13.0) [rounded corners={2pt}] to (0.0,1.0) to node[host label] {1} (9.0,1.0) [sharp corners] to (9.0,-13.0) -- cycle;
\path[host background] (14.5,-13.0) [rounded corners={2pt}] to (14.5,1.0) to node[host label] {2} (29.5,1.0) [sharp corners] to (29.5,-13.0) -- cycle;
\path[host background] (0.0,-6.5) [rounded corners={2pt}] to (0.0,-27.5) to (4.5,-27.5) [sharp corners] to (4.5,-31.5) to (19.0,-31.5) [rounded corners={2pt}] to (19.0,-27.5) to (29.5,-27.5) [sharp corners] to (29.5,-6.5) to (14.5,-6.5) [rounded corners={2pt}] to (14.5,-13.0) to (9.0,-13.0) [sharp corners] to (9.0,-6.5) -- cycle;
\path[host background] (33.5,-44.0) [rounded corners={2pt}] to (33.5,1.0) to node[host label] {3} (48.5,1.0) [sharp corners] to (48.5,-44.0) -- cycle;
\path[host background] (4.5,-22.25) [rounded corners={2pt}] to (4.5,-50.0) to (23.375,-50.0) [sharp corners] to (23.375,-54.0) to (29.375,-54.0) [rounded corners={2pt}] to (29.375,-50.0) to (48.5,-50.0) [sharp corners] to (48.5,-22.0) to (33.5,-22.0) [rounded corners={2pt}] to (33.5,-44.0) to (19.0,-44.0) [sharp corners] to (19.0,-22.25) -- cycle;
% unsampled hosts
% branches
\path[branch] (19.0,-9.0) to (19.0,-4.5);
\path[branch] (11.75,-20.25) [rounded corners={1.5pt}] to (4.5,-20.25) [sharp corners] to (4.5,-4.5);
\path[branch] (11.75,-20.25) [rounded corners={1.5pt}] to (25.0,-20.25) [sharp corners] to (25.0,-4.5);
\path[branch] (11.75,-27.5) to (11.75,-20.25);
\path[branch] (41.0,-27.5) to (41.0,-4.5);
\path[transfer branch] (41.0,-27.5) to (11.75,-27.5);
\path[branch] (41.0,-39.5) to (41.0,-27.5);
\path[invalid transfer branch] (41.0,-39.5) [out=180, in=0] to (19.0,-9.0);
% gene transfers
% events
\node[extant={black}{}] at (4.5,-4.5) {};
\node[extant={black}{}] at (19.0,-4.5) {};
\node[extant={black}{}] at (25.0,-4.5) {};
\node[speciation={black}{}] at (11.75,-20.25) {};
\node[extant={black}{}] at (41.0,-4.5) {};
\node[transfer={black}{}] at (41.0,-27.5) {};
\node[transfer={black}{}] at (41.0,-39.5) {};

    \end{scope}
    \draw[separator rule] (3.5, -.2) -- ++(0, 2.45);
    \begin{scope}[shift={(5, 0)}]
        \tikzset{
    x={(1pt,0)},
    y={(0,-1pt)},
    limit width/.style 2 args={
        execute at begin node={\begin{varwidth}{#1}#2},
        execute at end node={\end{varwidth}},
    },
    host background/.style={
        fill=host background color,
        draw=host background color,
        line width={.01pt},
    },
    unsampled host background/.style={
        preaction={fill=host background color},
        draw=host background color,
        pattern={Lines[angle=45, distance=3pt, line width=1.5pt]},
        pattern color=white,
    },
    host label/.style={
        font=\bfseries,
        midway,
        anchor=north,
        limit width={12em}{\centering},
    },
    branch/.style={
        draw={#1},
        line width={0.5pt},
    },
    loss branch/.style={
        branch,
        draw={black!40},
    },
    transfer branch/.style={
        branch={#1},
        densely dashed,
        -{Latex[length=0pt 8]},
    },
    invalid transfer branch/.style={
        transfer branch={red!80!black},
        out distance=20,
        in distance=20,
    },
    event/.style={
        draw={#1}, fill={host background color!50!white},
        font={\color{#1}},
        outer sep=0pt, inner xsep=0pt, inner ysep=2pt,
        limit width={6em}{\narrowragged},
        line width={0.5pt},
    },
    event/.default={black},
    extant/.code 2 args={
        \pgfkeysalso{
            circle, fill={#1},
            outer sep=0pt, inner sep=0pt,
            minimum size={3},
        }
        \def\cmp{#2}\ifx\cmp\empty\else
            \pgfkeysalso{
                label={[%
                    font={\color{#1}\vphantom{gb}},
                    fill=host background color!35!white,
                    rounded corners=2,
                    label distance=2pt,
                    inner xsep=2pt, inner ysep=1pt,
                    outer xsep=0pt, outer ysep=0pt,
                    limit width={6em}{\centering},
                ]below:#2},
            }
        \fi
    },
    extant/.default={black}{},
    unsampled/.style={
        circle, fill={#1},
        outer sep=0pt, inner sep=0pt,
        minimum size={3},
    },
    unsampled/.default={black!40},
    speciation/.style 2 args={
        event={#1}, rectangle, rounded corners,
        inner xsep=4pt,
        minimum width={8},
        minimum height={8},
        execute at begin node={#2},
    },
    duplication/.style 2 args={
        event={#1}, rectangle,
        inner xsep=4pt,
        minimum width={8},
        minimum height={8},
        execute at begin node={#2},
    },
    transfer/.code 2 args={
        \pgfkeysalso{
            event={#1},
        }
        \def\cmp{#2}\ifx\cmp\empty
            \pgfkeysalso{
                diamond,
                inner sep=1.5pt,
            }
        \else
            \pgfkeysalso{
                chamfered rectangle,
                chamfered rectangle sep={8 / 2.4},
                inner xsep=2pt,
                inner ysep=-1pt,
                minimum width={8},
                minimum height={8},
                execute at begin node={#2},
            }
        \fi
    },
    gain/.code 2 args={
        \pgfkeysalso{
            draw={#1}, fill=black,
            semicircle, inner sep=1pt,
            shape border rotate={0},
        }
        \def\cmp{#2}\ifx\cmp\empty\else
            \pgfkeysalso{
                label={[%
                    font={\scriptsize\vphantom{gb}},
                    label distance=2pt,
                    inner xsep=0pt, inner ysep=0pt,
                    outer xsep=0pt, outer ysep=0pt,
                    limit width={6em}{\narrowragged},
                ]right:#2},
            }
        \fi
    },
    loss/.code 2 args={
        \pgfkeysalso{
            draw={#1}, fill=black,
            semicircle, inner sep=1pt,
            shape border rotate={180+0},
        }
        \def\cmp{#2}\ifx\cmp\empty
            \pgfkeysalso{draw={black!40}, fill={black!40}}
        \else
            \pgfkeysalso{
                label={[%
                    font={\scriptsize\vphantom{gb}},
                    label distance=2pt,
                    inner xsep=0pt, inner ysep=0pt,
                    outer xsep=0pt, outer ysep=0pt,
                    limit width={6em}{\narrowragged},
                ]right:#2},
            }
        \fi
    },
}
% hosts
\path[host background] (0.0,-13.0) [rounded corners={2pt}] to (0.0,1.0) to node[host label] {1} (16.8722,1.0) [sharp corners] to (16.8722,-13.0) -- cycle;
\path[host background] (29.6222,-13.0) [rounded corners={2pt}] to (29.6222,1.0) to node[host label] {2} (44.6222,1.0) [sharp corners] to (44.6222,-13.0) -- cycle;
\path[host background] (0.0,-6.5) [rounded corners={2pt}] to (0.0,-39.0) to (12.3722,-39.0) [sharp corners] to (12.3722,-51.0) to (34.1222,-51.0) [rounded corners={2pt}] to (34.1222,-39.0) to (44.6222,-39.0) [sharp corners] to (44.6222,-6.5) to (29.6222,-6.5) [rounded corners={2pt}] to (29.6222,-13.0) to (16.8722,-13.0) [sharp corners] to (16.8722,-6.5) -- cycle;
\path[host background] (48.6222,-55.5) [rounded corners={2pt}] to (48.6222,1.0) to node[host label] {3} (63.6222,1.0) [sharp corners] to (63.6222,-55.5) -- cycle;
\path[host background] (12.3722,-32.0) [rounded corners={2pt}] to (12.3722,-61.5) to (38.4972,-61.5) [sharp corners] to (38.4972,-65.5) to (44.4972,-65.5) [rounded corners={2pt}] to (44.4972,-61.5) to (63.6222,-61.5) [sharp corners] to (63.6222,-27.75) to (48.6222,-27.75) [rounded corners={2pt}] to (48.6222,-55.5) to (34.1222,-55.5) [sharp corners] to (34.1222,-32.0) -- cycle;
% unsampled hosts
% branches
\path[branch] (19.6222,-20.25) [rounded corners={1.5pt}] to (34.1222,-20.25) [sharp corners] to (34.1222,-4.5);
\path[loss branch] (19.6222,-20.25) [rounded corners={1.5pt}] to (12.3722,-20.25) [sharp corners] to (12.3722,-4.4361);
\path[branch] (19.6222,-51.0) to (19.6222,-20.25);
\path[branch] (26.8722,-31.75) [rounded corners={1.5pt}] to (5.4361,-31.75) [sharp corners] to (5.4361,-4.5);
\path[branch] (26.8722,-31.75) [rounded corners={1.5pt}] to (40.1222,-31.75) [sharp corners] to (40.1222,-4.5);
\path[branch] (26.8722,-39.0) to (26.8722,-31.75);
\path[branch] (56.1222,-39.0) to (56.1222,-4.5);
\path[transfer branch] (56.1222,-39.0) to (26.8722,-39.0);
\path[branch] (56.1222,-51.0) to (56.1222,-39.0);
\path[transfer branch] (56.1222,-51.0) to (19.6222,-51.0);
% gene transfers
% events
\node[loss={black}{}] at (12.3722,-4.4361) {};
\node[extant={black}{}] at (5.4361,-4.5) {};
\node[extant={black}{}] at (34.1222,-4.5) {};
\node[extant={black}{}] at (40.1222,-4.5) {};
\node[speciation={black}{}] at (19.6222,-20.25) {};
\node[speciation={black}{}] at (26.8722,-31.75) {};
\node[extant={black}{}] at (56.1222,-4.5) {};
\node[transfer={black}{}] at (56.1222,-39.0) {};
\node[transfer={black}{}] at (56.1222,-51.0) {};

    \end{scope}
\end{tikzpicture}

    \caption[Examples of time inconsistencies]{
        Examples of time inconsistencies.
        \emph{(Left.)}
        Two alternative interpretations of the same reconciliation under the~\(\recmodelDTL\) model.
        The interpretation on the left is time-inconsistent whereas the one on the right is consistent.
        These diagrams are generated by the \emph{superrec2} software, presented in \cref{sec:logiciel:superrec2}, which automatically detects whether a reconciled tree is time-consistent.
        \emph{(Right.)}
        A reconciliation for which there exists no time-consistent interpretation.
    }
    \label{fig:reconciliation:time-cons}
\end{figure}

The \textsc{Bellman--Ford} algorithm makes it possible to check whether a given reconciled tree satisfies the conditions of \cref{def:reconciliation:time-cons} in \(\bigo((|\tnodes(\rectree)| + |\tnodes(S)|)^2)\) time~\parencite[section~22.4]{Cormen2022IntroductionAlgorithms}.
When the conditions are satisfied, this algorithm furthermore makes it possible to obtain a minimal total order satisfying the constraints and thus to help draw a diagram of the tree.

A reconciliation \(\rec\) is said to be \emph{acyclic} in a model \(\recmodel\) if there exists a time-consistent reconciled tree \(\rectree_{\!\recmodel}(\rec)\).
For example, the reconciliation \(\rec_3\) of \cref{fig:reconciliation:free-rec} is acyclic because of the interpretation given in \cref{fig:reconciliation:rec-trees}.
Some acyclic reconciliations, such as the one on the left of \cref{fig:reconciliation:time-cons}, have both consistent and inconsistent interpretations.

The following lemma, adapted from \textcite[theorem 1]{Tofigh2009InferringDuplicationsLateral}, makes it possible to check whether a reconciliation is acyclic in the \(\recmodelDTL\) model without enumerating all of its interpretations.

\begin{lemma}[Acyclicity of a reconciliation]\label{lem:reconciliation:rec-cons}
    A reconciliation \(\rec\) is acyclic in the \(\recmodelDTL\) model if and only if there exists a preorder \(\succeq_\sigma\) on \(\tnodes(S)\) such that:
    \begin{condition}
        \item \(s \succ_\sigma s\child\) for each branch \((s, s\child) \in \tedges(S)\); \label{lem:reconciliation:rec-cons:hosts}
        \item \(\tparent(\rec(v)) \succ_\sigma \rec(v'\child)\) for all branches \((v, v\child), (v', v'\child) \in \tedges(T)\) with \(v_c \tancs v_c'\) and \(\rec(v) \neq \troot(S)\).
            \label{lem:reconciliation:rec-cons:host-switch}
    \end{condition}
\end{lemma}
\begin{proof}
    Suppose that there exists an interpretation of \(\rec\) with a preorder \(\succeq_\tau\) satisfying the conditions of \cref{def:reconciliation:time-cons}.
    Observe that \cref{def:reconciliation:time-cons:hosts} of the definition is identical to \cref{lem:reconciliation:rec-cons:hosts} of the lemma, which is thus trivially satisfied.
    Moreover, let \((v, v\child), (v', v'\child) \in \tedges(T)\) be two branches such that \(v\child \tancs v'\child\) and \(\rec(v) \neq \troot(S)\).
    By \cref{def:reconciliation:time-cons:events,def:reconciliation:time-cons:events-hosts} of the definition, \(\tparent(\rec(v)) = \tparent(\exrec(v)) \succ_\tau v \succeq_\tau v\child\) and \(v'\child \succeq_\tau \exrec(v'\child) = \rec(v'\child)\) and therefore \(\tparent(\rec(v)) \succ_\tau \rec(v'\child)\).
    Thus, the restriction of~\(\succeq_\tau\) to the nodes of \(S\) gives a preorder satisfying the conditions of the lemma.

    Conversely, suppose that there exists a preorder \(\succeq_\sigma\) satisfying the conditions of the lemma and take a linear extension \(\succeq_\sigma'\) of it.
    We construct an interpretation \(\exrec\) of \(\rec\), not necessarily parsimonious, associated with a total order \(\succeq_\tau\) that completes \(\succeq_\sigma'\) and satisfies \cref{def:reconciliation:time-cons}.

    Following a prefix order on \(T\), for each host-switch branch \((v, v\child) \in \tedges(T)\) such that \(\rec(v) \nottcomp \rec(v\child)\), let us consider the time interval in which this switch can occur.
    This interval is bounded by \(b(v\child)\) and \(a(v\child)\), defined as
    \begin{align*}
        b(v\child) &= \min_{\succeq_\tau}\left(\{v' \mid v'\child \tancsneq v\child \text{ and } \exrec(v') \nottcomp \exrec(v'\child)\} \cup \{\tparent(\exrec(v))\}\right) \\
        a(v\child) &= \max_{\succeq_\tau}\left(\{\exrec(v'\child) \mid v\child \tancs v'\child \text{ and } \exrec(v') \nottcomp \exrec(v'\child)\} \cup \{\exrec(v)\}\right).
    \end{align*}
    Now note that \(b(v\child) \succ_\tau a(v\child)\).
    Let \((v^b, v^b\child) \in \tedges(T)\) and \((v^a, v^a\child) \in \tedges(T)\) be two branches such that \(v^b\child \tancsneq v\child\), \(v\child \tancs v^a\child\), \(\exrec(v^b) \nottcomp \exrec(v^b\child)\), and \(\exrec(v^a) \nottcomp \exrec(v^a\child)\).
    Observe that \(a(v^b\child) \succeq_\tau a(v\child)\), on the one hand because \(v^b\child \tancsneq v\child \tancs v'\child\), and on the other hand because \(\exrec(v''\child) \succeq_\tau \exrec(v)\), where \((v'', v''\child)\) is the most recent host-switch branch before \((v, v\child)\).
    Now, inductively, we can assume that \(v^b\) has~been placed such that \(v^b \succ_\tau a(v^b\child) \succeq_\tau a(v\child)\).
    By \cref{lem:reconciliation:rec-cons:host-switch} of the lemma, \(\tparent(\exrec(v)) \succ_\tau \exrec(v^a\child)\).
    Finally, \(\tparent(\exrec(v)) \succ_\tau \exrec(v)\) by \cref{lem:reconciliation:rec-cons:hosts} of the lemma.

    To interpret the host-switch branch \((v, v\child)\),
    it is subdivided by inserting two consecutive intermediate nodes \(v_0\) and \(v_1\).
    The node \(v_0\) is marked as a transfer event associated with the host \(\exrec(v_0) = \rec(v)\), leading on the one hand to a loss and on the other hand to the node~\(v_1\) resulting from the transfer.
    The non-empty interval between \(b(v\child)\) and \(a(v\child)\) makes it possible to date these nodes such that \(b(v\child) \succ_\tau v_0 =_\tau v_1 \succ_\tau a(v\child)\).
    This choice inductively satisfies \cref{def:reconciliation:time-cons:events} since \(v^b \succeq_\tau b(v\child) \succ_\tau v_0 =_\tau v_1\).
    \Cref{def:reconciliation:time-cons:events-hosts} of the definition is satisfied for~\(v_0\) since \(\tparent(\exrec(v)) \succeq_\tau b(v\child) \succ_\tau v_0 \succ_\tau a(v\child) \succeq_\tau \exrec(v)\).
    Finally, \(\lca\{\exrec(v), \exrec(v\child)\} \succeq_\tau \tparent(\exrec(v)) \succeq_\tau b(v\child)\), so there exists a host \(t\) such that \(\lca\{\exrec(v), \exrec(v\child)\} \tancsneq t \tancs \exrec(v\child)\) and \(\tparent(t) \succ_\tau v_1 \succeq_\tau t\).
    The node \(v_1\) is then associated with such a host \(t\), which satisfies \cref{def:reconciliation:time-cons:events-hosts} of the definition for~\(v_1\).

    Finally, the other nodes \(v \in \tinner(T)\) are marked as duplications and the branches \((v, v\child)\) such that \(\exrec(v) \tancsneq \exrec(v\child)\) are subdivided by inserting enough speciations and losses and by choosing suitable dates.
\end{proof}

For example, in the reconciliation \(\rec_5\) of \cref{fig:reconciliation:time-cons}, a first transfer branch from host~2 to host~3 is followed by another transfer from host~3 to host~\(\beta\).
These two transfers, in order to satisfy \cref{lem:reconciliation:rec-cons:host-switch} of the previous lemma, require that \(\beta >_\sigma \beta\), so there exists no interpretation of \(\rec_5\) that is time-consistent.

A reconciliation can, however, be cyclic without containing two transfers following one another.
For example, for the reconciliation \(\rec_4\) of \cref{fig:reconciliation:free-rec}, the transfer branch from host~4 to host~\(\beta\) requires that \(\gamma \succ_\sigma \beta\) whereas the one from host~1 to host~\(\gamma\) requires that \(\beta \succ_\sigma \gamma\).
Thus, the reconciliation \(\rec_4\) is also cyclic by the previous lemma.

A weaker version of \cref{lem:reconciliation:rec-cons:host-switch}, which did not consider pairs of transfers following one another, was initially stated by \textcite{Charleston1998JunglesNewSolution}.
As shown by \textcite{Tofigh2009InferringDuplicationsLateral}, this condition is incomplete, in particular because it does not identify reconciliation \(\rec_5\) as cyclic.

\subsection{Complexity of \(\recmodel\)-parsimony problems}

As suggested by the following two statements, the space of strict reconciliations is more easily explorable and more restricted than the general space of reconciliations.

\begin{remark}
    All strict reconciliations are acyclic.
\end{remark}

\begin{lemma}\label{lemma:strict-rec-space}
    If \(\rec\) is a strict reconciliation of \(T\) with \(S\) then, for every internal node \(v \in \tnodes(T)\),
    \(\rec(v) \tancs \lca\left\{ \llmap(l) \suchthat l \in \tleavs(\subtree{T}{v}) \right\}\).
\end{lemma}
\vspace{-1em}
\begin{proof}
    Let \(\rec\) be a reconciliation of \(T\) with \(S\) and suppose that an internal node \(v \in \tnodes(T)\) is such that
    \(\rec(v) \nottancs \lca\left\{ \llmap(l) \suchthat l \in \tleavs(\subtree{T}{v}) \right\}\).
    There therefore exists \(w \in \tleavs(\subtree{T}{v})\) such that \(\rec(v) \nottancs \llmap(w)\)
    and therefore such that \(\rec(v) \not\tancs \rec(w)\).
    Hence \(\rec\) is not strict.
\end{proof}

\vspace{-.5em}
The strict reconciliation that associates each internal node of \(T\) with the most recent possible host in~\(S\) is denoted \(\reclca\) and defined by \(
    \reclca(v) = \lca\left\{ \llmap(l) \suchthat l \in \tleavs(\subtree{T}{v}) \right\}
\).

Since the \(\recmodelDL\) model excludes any possibility of host switch, the optimal reconciliations lie within the space of strict reconciliations (see lemma~\labelcref{lemma:strict-rec-space}).
The reconciliation~\(\reclca\) is always a solution under this model.
Moreover, provided that \(\cost(\evtL) > 0\), it is the unique solution~\parencite{Chauve2009NewPerspectivesGene}.
Computing~\(\reclca\) only requires linear time and space in the size of \(T\)~\parencite{Zhang1997MirkinMuchnikSmith}.
The \(\recmodelDL\)-parsimony problem is therefore polynomial.

Unlike the \(\recmodelDL\) model, there may exist several optimal reconciliations under the \(\recmodelDTL\) model and the reconciliation \(\reclca\) is not necessarily one of them.
Moreover, a natural constraint for this problem is to only take acyclic reconciliations into consideration.
However, this constraint makes the associated decision problem NP-complete, as shown by \textcite{Tofigh2011SimultaneousIdentificationDuplications} through a reduction from the maximum acyclic subgraph problem.

To circumvent this difficulty, in what follows, we will allow cyclic reconciliations.
We will also evaluate the cost of a reconciliation as that of its least costly interpretation, even if the latter is inconsistent.
Thus, the reconciliation \(\rec_4\) of \cref{fig:reconciliation:free-rec} becomes acceptable and the reconciliation of \cref{fig:reconciliation:time-cons} costs only two transfer events, even though a loss is necessary to make it a time-consistent interpretation.

These modifications give the problem an optimal substructure: a solution for the subtree \(\subtree{T}{v}\) is always an assembly of two solutions for \(\subtree{T}{v\lchild}\) and \(\subtree{T}{v\rchild}\).
\Cref{alg:reconciliation:dtl-parcimonie} exploits this property to solve the problem efficiently.

\newcommand{\fndtl}{\mathrm{cost}}
\newcommand{\varspe}{\mathrm{Spe_{LR}}}
\newcommand{\varsperl}{\mathrm{Spe_{RL}}}
\newcommand{\vardup}{\mathrm{Dup}}
\newcommand{\varhgt}{\mathrm{Tra_L}}
\newcommand{\varhgtr}{\mathrm{Tra_R}}
\newcommand{\varhgtlr}{\mathrm{Tra_{LR}}}

\begin{algorithm}
    \capstart
    \begin{algorithmic}[1]\setstretch{1.25}
        \Require{%
            \(T\) and \(S\), two phylogenies;\\
            \(\llmap : \tleavs(T) \onto \tleavs(S)\), an association between their leaves;\\
            \(\cost : \recmodelDTL \onto \realpos\), a cost for the \(\recmodelDTL\) events.
        }
        \Ensure{Minimum cost of a reconciliation of \(T\) with \(S\) under the \(\recmodelDTL\) model, including cyclic reconciliations.}
        \Function{cost}{\(v\), \(\sigma\)}
            \If{\(v \in \tleavs(T)\)}
                \Return \(0\) {\algorithmicif} \(\sigma = \llmap(v)\)\textbf{,} {\algorithmicelse} \(\infty\) \;
            \EndIf
            \State \(\displaystyle\varspe \gets \min_{\sigma' \tdesc \sigma\lchild} \{
                \fndtl(v\lchild, \sigma') + \cost(\evtL) \cdot \tdist(\sigma\lchild, \sigma')
            \} + \min_{\sigma'' \tdesc \sigma\rchild} \{
                \fndtl(v\rchild\kern1pt, \sigma'') + \cost(\evtL) \cdot \tdist(\sigma\rchild\kern1pt, \sigma'')
            \}\)
            \State \(\displaystyle\varsperl \gets \min_{\sigma' \tdesc \sigma\rchild} \{
                \fndtl(v\lchild, \sigma') + \cost(\evtL) \cdot \tdist(\sigma\rchild\kern1pt, \sigma')
            \} + \min_{\sigma'' \tdesc \sigma\lchild} \{
                \fndtl(v\rchild\kern1pt, \sigma'') + \cost(\evtL) \cdot \tdist(\sigma\lchild, \sigma'')
            \}\)
            \State \(\displaystyle\vardup \gets \cost(\evtD) + \min_{\sigma' \tdesc \sigma} \{
                \fndtl(v\lchild, \sigma') + \cost(\evtL) \cdot \tdist(\sigma, \sigma')
            \} + \min_{\sigma'' \tdesc \sigma} \{
                \fndtl(v\rchild\kern1pt, \sigma'') + \cost(\evtL) \cdot \tdist(\sigma, \sigma'')
            \}\)
            \State \(\displaystyle\varhgt \gets \cost(\evtT) + \min_{\sigma' \nottcomp \sigma} \{
                \fndtl(v\lchild, \sigma')
            \} + \min_{\sigma'' \tdesc \sigma} \{
                \fndtl(v\rchild\kern1pt, \sigma'') + \cost(\evtL) \cdot \tdist(\sigma, \sigma'')
            \}\)
            \State \(\displaystyle\varhgtr \gets \cost(\evtT) + \min_{\sigma' \tdesc \sigma} \{
                \fndtl(v\lchild, \sigma') + \cost(\evtL) \cdot \tdist(\sigma, \sigma')
            \} + \min_{\sigma'' \nottcomp \sigma} \{
                \fndtl(v\rchild\kern1pt, \sigma'')
            \}\)
            \State \Return \(\min\{\varspe, \varsperl, \vardup, \varhgt, \varhgtr\}\)
        \EndFunction
        \State \Return \(\min_{\sigma \in \tnodes(S)}\Call{cost}{\troot(T), \sigma}\)
    \end{algorithmic}
    \caption{
        Recursive algorithm for the \(\recmodelDTL\)-parsimony problem.
        \label{alg:reconciliation:dtl-parcimonie}
    }
\end{algorithm}

The function \(\Call{cost}{v, \sigma}\) computes the minimum cost of a reconciliation of \(\subtree{T}{v}\) with \(S\) such that \(v\) is associated with the given host \(\sigma\).
To do so, recursive calls are used to obtain the minimum cost of reconciliations of \(\subtree{T}{v\lchild}\) and \(\subtree{T}{v\rchild}\) with \(S\) when \(v\lchild\) and \(v\rchild\) are associated with different hosts \(\sigma'\) and \(\sigma''\) of \(S\).
The assembly of sub-reconciliations in the algorithm is based on the following observations:
\begin{compactitemize}
    \item \(\evt(v) = \evtS\) only if \(\sigma\) is an internal node and \(\sigma'\) and \(\sigma''\) are each in one of the two subtrees \(\subtree{S}{\sigma\lchild}\) and \(\subtree{S}{\sigma\rchild}\) (variables \(\varspe\) and \(\varsperl\));
    \item \(\evt(v) = \evtD\) only if \(\sigma'\) and \(\sigma''\) descend from \(\sigma\) (variable \(\vardup\));
    \item \(\evt(v) = \evtT\) only if \(\sigma'\) or \(\sigma''\) are separated from \(\sigma\) (variables \(\varhgt\) and \(\varhgtr\));
    \item except for transfer branches, loss events must be inserted for each host traversed between \(\sigma\), \(\sigma'\) and \(\sigma''\) in order to comply with the model (terms in \(\cost(\evtL) \cdot \hdots\)).
\end{compactitemize}

\Cref{alg:reconciliation:dtl-parcimonie} can be transformed to use a dynamic programming table that stores the values of \(\fndtl(v, \sigma)\), thereby avoiding repeating the same computation several times.
Letting \(n = \lvert\tnodes(T)\rvert\) and \(m = \lvert\tnodes(S)\rvert\), the algorithm thus transformed has a time complexity of \(\bigo(nm^2)\) and a space complexity of \(\bigo(nm)\)~\parencite{Tofigh2011SimultaneousIdentificationDuplications}.
The time complexity can be reduced to \(\bigo(nm)\) without modifying the space complexity~\parencite{Bansal2012EfficientAlgorithmsReconciliation}.
The \(\recmodelDTL\)-parsimony problem is therefore polynomial.

\Cref{alg:reconciliation:dtl-parcimonie} can be extended to the \(\recmodelDTLx\) model without changing its complexity.
To do so, it suffices to add, for each of lines~3 to~7, a case in which the terms in \(\cost(\evtL) \cdot \hdots\) are replaced by the invocation of a \(\evtTX\) event.

\subsection{Applications of reconciliation methods}

Bilateral reconciliation allows the study of the joint evolution of evolutionary units and their hosts.
The algorithms described in the previous section make it possible to obtain information about these evolutionary relationships from the respective phylogenies of the units and of the hosts.
Under some evolutionary models, such as the~\(\recmodelDTL\) model, several distinct solutions may, however, be optimal in the sense of the parsimony criterion.
Moreover, the execution of these algorithms depends on the costs set for each event type, and varying these costs modifies the optimal solutions.
It is therefore not sufficient to analyse a single solution returned by a single parameterization of an algorithm.
For a meaningful analysis, the solution space must be explored to quantify the extent to which these solutions agree on the different parts of the inferred scenario~\parencite[see][]{Doyon2009SpaceGeneSpeciesTrees,Weiner2021ImprovedDuplicationTransfer}.
We will address these questions from the angle of algebraic dynamic programming in \cref{chap:adp,chap:adp-rec}.

The phylogenies expected as input to the algorithms are sometimes known with limited precision, or are even unknown.
Reconciliation can provide a criterion for improving these phylogenies.
For example, reconciliation can be used to infer an unknown species tree from a set of trees of several gene families~\parencite[see][p.~72--73]{Guigo1996ReconstructionAncientMolecular,Page1994MapsTreesCladistic}.
The underlying problem being NP-complete~\parencite{Ma2000GeneTreesSpecies}, this application of reconciliation relies on heuristics.
In some cases, particularly when there are uncertainties in its construction, the phylogeny of the units may contain polytomies, that is, internal nodes having more than two descendants.
Reconciliation can then be used as a criterion to \emph{resolve} these polytomies by searching for the binary tree that yields a lower reconciliation score~\parencite{Lafond2016EfficientNonBinary}.

%! TeX root = ./main.tex
% vim: set spelllang=en:
\section{Multiple reconciliation}

Some hosts carry several families of units.
To reconstruct the respective histories of these families, bilateral reconciliation methods could be applied to them individually.
This approach would, however, assume that the units under study evolve independently of one another, which is in particular not the case for some genes.
Genes are indeed organized linearly along the genome of their species and affected by genome-scale events.
These events modify the organization and the composition of the genome, acting simultaneously on several of its genes.

The evolution of linear structures, shaped by rearrangements and duplications, is mainly studied in the context of comparative genomics.
The \emph{genome phylogeny} problem consists in reconstructing the most plausible phylogenetic tree for a set of genomes, by associating each internal node with an ancestral genome.
This problem comes in a \emph{\enquote{large}} version, where the topology of the tree and the ancestral genomes are inferred simultaneously, and a \emph{\enquote{small}} version, in which the topology is fixed and only the ancestral genomes remain to be determined~\parencite[for a review of these problems, see][]{ElMabrouk2012AnalysisGeneOrder}.

To address these problems, \textcite{Sankoff2000DuplicationRearrangementReconciliation} proposed combining the methods developed to analyse evolution at the scale of gene families with those for the scale of genomes.
The information recovered by reconciliation, in particular orthology relationships, could indeed guide the study of the evolution of genomes containing several copies of the same gene family.
Conversely, integrating the genomic scale into reconciliation would make it possible to take into account the interdependence of the evolution of genes.
The research objective is thus the development of methods capable of reconstructing the joint history of several gene families~\parencite[see also][]{Chauve2013DuplicationRearrangementReconciliation}.

To do so, an intermediate level can be integrated into the bilateral reconciliation model, placed between the scale of genes and that of species, in order to model evolution at the genomic scale.
In what follows, we compare two existing approaches based on this idea according to their modelling capacity and their algorithmic complexity.

\subsection{Adjacency reconciliation}

The linear organization of a set of units can be represented by a set of adjacencies between pairs of units.
For example, the genome \(\{(a, b, c), (d, e)\}\) made up of two chromosomes and five genes corresponds to the three adjacencies \(\{\{a, b\}, \{b, c\}, \{d, e\}\}\).
This representation notably makes it easier to model rearrangement events, since it suffices to act on the adjacencies at the border of the segment moved by the event.
The DeCo approach (for \enquote{Detection of Coevolution}), introduced by \textcite{Berard2012EvolutionGeneNeighborhoods}, extends reconciliation to integrate the evolution of adjacencies between pairs of dependent units.

Suppose that a forest \(\gtrees\) of phylogenies of dependent units and the phylogeny \(S\) of their hosts are known.
An adjacency \(\{a, b\}\) between two internal nodes of \(\gtrees\) is an assertion stating that the two divergences occurred simultaneously in the same host and that the copy of \(a\) is an immediate neighbour of \(b\) in the structure of this host.

\begin{definition}[Adjacencies]\label{definition:adj}
    An \emph{adjacency} on a set \(L\) is a pair \(\{a, b\} \subseteq L\) such that \(a \neq b\).
    A map \(f : L \onto M\) \emph{respects} such an adjacency if \(f(a) = f(b)\).
    The set of all possible adjacencies respected by~\(f\) is denoted~\(\alladjs(f)\).
\end{definition}

Suppose that the set of adjacencies between the leaves of \(\gtrees\) is also known.
An adjacency reconciliation of \(\gtrees\) with \(S\) (see \cref{fig:reconciliation:adj-rec}) is a reconciliation of each tree of \(\gtrees\) with \(S\) together with a set of ancestral adjacencies leading to the adjacencies between the leaves.

\begin{definition}[Adjacency reconciliations]\label{definition:adjrec}
    Let \(\gtrees\) be a phylogenetic forest and \(S\) a phylogenetic tree.
    Let \(\llmap : \tleavs(\gtrees) \onto \tleavs(S)\) be a map and \(\lladjset \subseteq \alladjs(\llmap)\) a set of adjacencies.
    An \emph{adjacency reconciliation} of \(\gtrees\) with \(S\) is a pair \(\tup{\rec, \adjset}\),
    where \(\rec\) is a map extending~\(\llmap\) that is the union of individual reconciliations of each tree of \(\gtrees\) with \(S\),
    and \(\adjset\) is a set of ancestral adjacencies such that \(\lladjset \subseteq \adjset \subseteq \alladjs(\rec)\).
\end{definition}

\begin{figure}[htb]
    \capstart
    \centering
    %! TeX root = ../main.tex
% vim: set spelllang=fr:
\begin{tikzpicture}
    %% Arbre T_a
    \begin{scope}[color1]
        \graph[
            phylogenetic tree layout,
            sibling distance=.2cm,
            evolutionary unit length=0.44cm,
            nodes={font=\small},
        ]{
            ta l4s1 / [coordinate] -- {
                ta l3s2 / [coordinate] -- {
                    ta l2s2 / [coordinate] -- {
                        ta l1s6 / \(\explar{a}{3}{2}\),
                        ta l1s5 / \(\explar{a}{2}{2}\),
                    },
                    ta l1s4 / \(\explar{a}{1}{2}\) [> length=2],
                },
                ta l3s1 / [coordinate] -- {
                    ta l2s1 / [coordinate] -- {
                        ta l1s3 / \(\explar{a}{3}{1}\),
                        ta l1s2 / \(\explar{a}{2}{1}\),
                    },
                    ta l1s1 / \(\explar{a}{1}{1}\) [> length=2],
                },
            },
        };
        \node at (0, -2.3) {\(T_a\)};
    \end{scope}

    %% Arbre T_b
    \begin{scope}[shift={(8, 0)}, color2]
        \graph[
            phylogenetic tree layout,
            sibling distance=.2cm,
            evolutionary unit length=0.44cm,
            nodes={font=\small},
        ]{
            tb l4s1 / [coordinate] -- {
                tb l3s2 / [coordinate] -- {
                    tb l2s2 / [coordinate] -- {
                        tb l1s6 / \(\explar{b}{3}{2}\),
                        tb l1s5 / \(\explar{b}{2}{1}\),
                    },
                    tb l1s4 / \(\explar{b}{1}{3}\) [> length=2],
                },
                tb l3s1 / [coordinate] -- {
                    tb l1s1 / \(\explar{b}{3}{1}\) [> length=2],
                    tb l2s1 / [coordinate] -- {
                        tb l1s3 / \(\explar{b}{1}{2}\),
                        tb l1s2 / \(\explar{b}{1}{1}\),
                    },
                },
            },
        };
        \node at (0, -2.3) {\(T_b\)};
    \end{scope}
\end{tikzpicture}

\vspace{1em}
\begin{tikzpicture}
    %% Arbre T_c
    \begin{scope}[color3]
        \graph[
            phylogenetic tree layout,
            sibling distance=.2cm,
            nodes={font=\small},
        ]{
            tc l3s1 / [coordinate] -- {
                tc l2s1 / [coordinate] -- {
                    tc l1s3 / \(\explar{c}{3}{1}\),
                    tc l1s2 / \(\explar{c}{2}{1}\),
                },
                tc l1s1 / \(\explar{c}{1}{1}\) [> length=2],
            },
        };
        \coordinate (labels) at (0, -2.3);
        \node at (labels-|tc l1s2) {\(T_c\)};
    \end{scope}

    %% Associations
    \node[anchor=north] at (4.25, 0) {\small\begin{tabular}{@{}l@{}l@{}}
        \(A = \{\)
        & \(\{{\color{color1}\explar{a}{1}{1}}, {\color{color2}\explar{b}{1}{1}}\},
            \{{\color{color2}\explar{b}{1}{1}}, {\color{color2}\explar{b}{1}{2}}\},
            \{{\color{color1}\explar{a}{1}{2}}, {\color{color2}\explar{b}{1}{3}}\},\) \\[3pt]
        & \(\{{\color{color1}\explar{a}{2}{2}}, {\color{color2}\explar{b}{2}{1}}\},
            \{{\color{color2}\explar{b}{2}{1}}, {\color{color3}\explar{c}{2}{1}}\},
            \{{\color{color1}\explar{a}{3}{1}}, {\color{color2}\explar{b}{3}{1}}\},\) \\[3pt]
        & \(\{{\color{color1}\explar{a}{3}{2}}, {\color{color2}\explar{b}{3}{2}}\},
            \{{\color{color2}\explar{b}{3}{2}}, {\color{color3}\explar{c}{3}{1}}\}\}\)
    \end{tabular}};

    %% Arbre S
    \begin{scope}[shift={(8, 0)}]
        \graph[
            phylogenetic tree layout,
            sibling distance=.2cm,
            nodes={font=\small\bfseries},
        ]{
            s l3s1 / [coordinate, label={[host]right:α}] -- {
                s l2s1 / [coordinate, label={[host]right:β}] -- {
                    s l1s3 / 3,
                    s l1s2 / 2,
                },
                s l1s1 / 1 [> length=2],
            },
        };
        \coordinate (labels) at (0, -2.3);
        \node at (labels-|s l1s2) {\(S\)};
    \end{scope}
\end{tikzpicture}

\vspace{1.5em}
\begin{tikzpicture}[x=.8cm, y=1cm]
    \draw[separator rule] (0, 3.95) -- ++(14, 0);

    %% Arbre T_a
    \begin{scope}[color1]
        \graph[no placement, nodes={font=\small}]{
            ta l4s1 / [graph node, x=5.5, y=3, label={[host, black, yshift=2pt]left:α}] -- {
                ta l3s1 / [graph node, x=2, y=2, label={[host, black]left:α}] -- {
                    ta l1s1 / \(\explar{a}{1}{1}\) [x=0],
                    ta l2s1 / [graph node, x=3.33, y=1, label={[host, black, yshift=-2pt]left:β}] -- {
                        ta l1s2 / \(\explar{a}{2}{1}\) [x=3],
                        ta l1s3 / \(\explar{a}{3}{1}\) [x=4],
                    },
                },
                ta l3s2 / [graph node, x=9, y=2] -- {
                    ta l1s4 / \(\explar{a}{1}{2}\) [x=6],
                    ta l2s2 / [graph node, x=10.5, y=1] -- {
                        ta l1s5 / \(\explar{a}{2}{2}\) [x=9],
                        ta l1s6 / \(\explar{a}{3}{2}\) [x=12],
                    },
                },
            },
        };
    \end{scope}

    %% Arbre T_b
    \begin{scope}[color2]
        \graph[no placement, nodes={font=\small}]{
            tb l4s1 / [graph node, x=6.5, y=3] -- {
                tb l3s1 / [graph node, x=3, y=2] -- {
                    tb l1s1 / \(\explar{b}{3}{1}\) [x=5],
                    tb l2s1 / [graph node, x=1.66, y=1, label={[host, black]right:1}] -- {
                        tb l1s2 / \(\explar{b}{1}{1}\) [x=1],
                        tb l1s3 / \(\explar{b}{1}{2}\) [x=2],
                    },
                },
                tb l3s2 / [graph node, x=10, y=2, label={[host, black]right:α}] -- {
                    tb l1s4 / \(\explar{b}{1}{3}\) [x=7],
                    tb l2s2 / [graph node, x=11.5, y=1] -- {
                        tb l1s5 / \(\explar{b}{2}{1}\) [x=10],
                        tb l1s6 / \(\explar{b}{3}{2}\) [x=13],
                    },
                },
            },
        };
    \end{scope}

    %% Arbre T_c
    \begin{scope}[color3]
        \graph[no placement, nodes={font=\small}]{
            tc l3s1 / [graph node, x=11, y=2, label={[host, black]right:α}] -- {
                tc l1s1 / \(\explar{c}{1}{1}\) [x=8],
                tc l2s1 / [graph node, x=12.5, y=1, label={[host, black, yshift=2pt]right:β}] -- {
                    tc l1s2 / \(\explar{c}{2}{1}\) [x=11],
                    tc l1s3 / \(\explar{c}{3}{1}\) [x=14],
                },
            },
        };
    \end{scope}

    %% Adjacences
    \begin{pgfonlayer}{hostedges}
        \path[adjacency] (ta l4s1) edge (tb l4s1);
        \path[adjacency] (ta l3s1) edge (tb l3s1);
        \path[adjacency] (ta l3s2) edge (tb l3s2);
        \path[adjacency] (ta l2s2) edge (tb l2s2);
        \path[adjacency] (tb l2s2) edge (tc l2s1);

        \path[adjacency] (ta l1s1) edge (tb l1s2);
        \path[adjacency] (tb l1s2) edge (tb l1s3);
        \path[adjacency] (ta l1s3) edge (tb l1s1);
        \path[adjacency] (ta l1s4) edge (tb l1s4);
        \path[adjacency] (ta l1s5) edge (tb l1s5);
        \path[adjacency] (tb l1s5) edge (tc l1s2);
        \path[adjacency] (ta l1s6) edge (tb l1s6);
        \path[adjacency] (tb l1s6) edge (tc l1s3);
    \end{pgfonlayer}
\end{tikzpicture}

    \caption[Instance of and solution to an adjacency reconciliation problem]{
        \emph{(Top.)}
        Forest of phylogenetic trees \(\gtrees = T_a \sqcup T_b \sqcup T_c\)
        and phylogeny \(S\).
        Each leaf of \(\gtrees\) is annotated with its host as a subscript
        and a copy number as a superscript, which makes it possible to distinguish it in the set of adjacencies \(\lladjset\).
        \emph{(Bottom.)}
        Adjacency reconciliation of \(\gtrees\) with \(S\).
        Each pair of adjacent nodes is connected by a dotted line.
        The reconciliation of the trees with \(S\) is indicated by labels.
    }
    \label{fig:reconciliation:adj-rec}
\end{figure}

A set of ancestral adjacencies implies temporal constraints that must be compatible with those imposed by the topology of \(\gtrees\).
Thus, a node cannot be adjacent to one of its descendants or ancestors and, if two nodes are adjacent, then an ancestor of one cannot be adjacent to a descendant of the other.

\vspace{1ex}
\begin{definition}[Compatibility]\label{definition:adjrec-adj-compat}
    Let \(\gtrees\) be a phylogenetic forest, \(S\) a phylogenetic tree, and \(\adjset\) a set of adjacencies of \(\gtrees\).
    The adjacencies of \(\adjset\) are \emph{compatible} if there exists a preorder~\(\succeq_\alpha\) on the nodes of \(\tnodes(\gtrees)\) such that:
    \begin{condition}
        \item \(v \succ_\alpha v\child\) for each branch \((v, v\child) \in \tedges(\gtrees)\);
        \item \(v \succeq_\alpha w\) and \(w \succeq_\alpha v\) for each adjacency \(\{v, w\} \in \adjset\). \qedhere
    \end{condition}
\end{definition}

Moreover, for them to correspond to a known biological structure, these adjacencies must be able to be assembled into segments or cycles without forming branchings.

\begin{definition}[Linearity]\label{definition:adjrec-lin}
    Let \(\gtrees\) be a phylogenetic forest, \(S\) a phylogenetic tree, and \(\tup{\rec, \adjset}\) an adjacency reconciliation of \(\gtrees\) with \(S\).
    The adjacencies of \(\adjset\) are \emph{linear} if each \(v \in \tnodes(\gtrees)\) appears in at most two adjacencies.
\end{definition}

As for bilateral reconciliations, choosing an evolutionary model makes it possible to interpret an adjacency reconciliation.
Such a model specifies not only the events that individually affect the dependent units, but also those that affect their adjacencies.
An interpretation is represented by an adjacency forest.

\begin{definition}[Adjacency forests]\label{definition:adjrectree}
    Let \(\gtrees\) be a phylogenetic forest, \(S\) a phylogenetic tree, \(\tup{\rec, \adjset}\) an adjacency reconciliation of \(\gtrees\) with \(S\), and \(\tup{\recmodel\rmunit, \recmodel\rmadj}\) two sets of event types.
    Let \(\rectree(\rec)\) be a reconciled forest under the model \(\recmodel\rmunit\) containing the maps \(\exrec\) and \(\evt\).
    An \emph{adjacency forest} \(\adjtree\) is composed of nodes
        labelled by adjacencies
        \(\exadj : \tnodes(\adjtree) \onto \alladjs(\exrec)\)
        and by events
        \(\evt\rmadj : \tnodes(\adjtree) \onto \recmodel\rmadj \cup \{\evtE, \evtB\}\),
        such that:
    \begin{condition}
        \item only the leaves of \(\adjtree\) can be labelled by the events \(\evtE\) or \(\evtB\), where \(\evtE\) labels the extant adjacencies and \(\evtB\) the adjacency breaks;
        \item \(\lladjset = \{\exadj(v) \suchthat v \in \tnodes(\adjtree)
            \text{ and } \evt(v) = \evtE \}\);
        \item \(\adjset \subseteq \{\exadj(v) \suchthat v \in \tnodes(\adjtree)
            \text{ and } \evt(v) \neq \evtB \}\). \qedhere
    \end{condition}
\end{definition}

Each root of an adjacency forest corresponds to the appearance of a new adjacency between two units\footnote{The model of \textcite{Berard2012EvolutionGeneNeighborhoods} makes it possible, under certain conditions, not to count some roots of the forest as new adjacencies but rather as tandem duplications.}, which then propagates according to the events of \(\recmodel\rmadj\).
This adjacency gain event is denoted~\(\evtG\).
Consequently, a forest without any branch corresponds to a scenario in which all adjacencies appeared spontaneously, whereas a forest made up of a single tree corresponds to the emergence of all adjacencies from a single ancestor.
The cost associated with an adjacency forest is the sum of the costs of its events, including as many gain events as it has roots.

\begin{problem}{\(\tup{\recmodel\rmunit, \recmodel\rmadj}\)-parsimony}
    \probinput{Phylogenies \(\gtrees\) and \(S\);
               leaf association \(\llmap : \tleavs(\gtrees) \onto \tleavs(S)\);
               adjacencies between leaves \(\lladjset \subseteq \alladjs(\llmap)\);
               costs \(\cost : \recmodel\rmunit \cup \recmodel\rmadj \cup \{\evtG\} \onto \nat\).}
    \proboutput{Minimum cost of an acyclic reconciled forest of \(\gtrees\) with \(S\) under the model \(\recmodel\rmunit\) and of an associated adjacency forest under the model \(\recmodel\rmadj\) for the costs \(\cost\).}
\end{problem}

A simple model for the evolution of adjacencies, based on duplication and loss events, was proposed by \textcite{Berard2012EvolutionGeneNeighborhoods}.
In this model, \(\recmodelDL\rmunit = \{\evtD, \evtL, \evtS\}\), \(\recmodelDL\rmadj = \{\evtDD, \evtD, \evtLL, \evtSS\}\), and each node \(v \in \tnodes(\adjtree)\) of a reconciled forest is such that:
\begin{compactitemize}
    \item if \(\evt\rmadj(v) = \evtDD\) (resp. \(\evtSS\)), then
        \(\exadj(v) = \{x, y\}\),
        \(\evt(x) = \evt(y) = \evtD\) (resp. \(\evtS\)), and
        \(\{\exadj(v\lchild), \exadj(v\rchild)\} = \{\{x\lchild, y\lchild\}, \{x\rchild, y\rchild\}\}\)
        or \(\{\exadj(v\lchild), \exadj(v\rchild)\} = \{\{x\lchild, y\rchild\}, \{x\rchild, y\lchild\}\}\);
    \item if \(\evt\rmadj(v) = \evtLL\), then \(v\) is a leaf,
        \(\exadj(v) = \{x, y\}\),
        and \(\evt(x) = \evt(y) = \evtL\);
    \item if \(\evt\rmadj(v) = \evtD\), then \(v\) is unary,
        \(\exadj(v) = \{x, y\}\),
        \(\evt(x) = \evtD\),
        and \(\exadj(v\child) = \{x', y\}\),
        where \(x' \in \children(x)\).
\end{compactitemize}

\Cref{fig:reconciliation:adj-forest} gives an example of a reconciled forest and of an adjacency forest interpreting the adjacency reconciliation of \cref{fig:reconciliation:adj-rec} under the model \(\tup{\recmodelDL\rmunit, \recmodelDL\rmadj}\).

\begin{figure}[htb]
    \capstart
    \centering
    %! TeX root = ../main.tex
% vim: set spelllang=fr:
\begin{tikzpicture}[
    label distance=-2pt,
    reference/.style={font=\footnotesize, opacity=.5},
]
    %% Réconciliation T_a
    \begin{scope}[x=.5cm, y=.8cm]
        \graph[no placement, nodes={dependent}]{
            tal4s1 / [duplication, x=3.5, y=2.75, label={[reference]left:\(\bar{a}_1\)}] -- {
                tal3s1 / [speciation, x=3, y=2, label={[reference]left:\(\bar{a}_2\)}] -- {
                    tal2s1 / [graph node, x=1.5, y=1, label={[reference, yshift=1pt]left:\(\bar{a}_4\)}] --
                        tal1s1 / \(\explar{a}{1}{1}\) [x=0, y=0, >shorten >=-3pt],
                    tal2s2 / [speciation, x=5.75, y=1.1, label={[reference]left:\(\bar{a}_5\)}] -- {
                        tal1s2 / \(\explar{a}{2}{1}\) [x=4, y=0, >shorten >=-3pt],
                        tal1s3 / \(\explar{a}{3}{1}\) [x=7.5, y=0, >shorten >=-3pt],
                    },
                },
                tal3s2 / [speciation, x=4, y=2, label={[reference]right:\(\bar{a}_3\)}] -- {
                    tal1s4 / \(\explar{a}{1}{2}\) [x=1, y=0, >shorten >=-3pt],
                    tal2s3 / [speciation, x=6.75, y=1.1, label={[reference]right:\(\bar{a}_6\)}] -- {
                        tal1s5 / \(\explar{a}{2}{2}\) [x=5, y=0, >shorten >=-3pt],
                        tal1s6 / \(\explar{a}{3}{2}\) [x=8.5, y=0, >shorten >=-3pt],
                    },
                },
            },
        };
    \end{scope}

    %% Réconciliation T_b
    \begin{scope}[x=.5cm, y=.8cm, shift={(12, 0)}]
        \graph[no placement, nodes={dependent}]{
            tbl4s1 / [duplication, x=3.5, y=2.75, label={[reference]left:\(\bar{b}_1\)}] -- {
                tbl3s1 / [speciation, x=3, y=2, label={[reference, yshift=1pt]left:\(\bar{b}_2\)}] -- {
                    tbl2s1 / [duplication, x=1.5, y=1, label={[reference, yshift=1pt]left:\(\bar{b}_4\)}] -- {
                        tbl1s1 / \(\explar{b}{1}{1}\) [x=-.5, y=0, >shorten >=-3pt],
                        tbl1s2 / \(\explar{b}{1}{2}\) [x=.5, y=0, >shorten >=-3pt],
                    },
                    tbl2s2 / [speciation, x=5.75, y=1.1, label={[reference, xshift=-1pt]left:\(\bar{b}_5\)}] -- {
                        tbl1s3 / [loss, >loss branch, rotate=-45, x=4, y=0],
                        tbl1s4 / \(\explar{b}{3}{1}\) [x=7.5, y=0, >shorten >=-3pt],
                    },
                },
                tbl3s2 / [speciation, x=4, y=2, label={[reference, yshift=1pt]right:\(\bar{b}_3\)}] -- {
                    tbl1s5 / \(\explar{b}{1}{3}\) [x=1.5, y=0, >shorten >=-3pt],
                    tbl2s3 / [speciation, x=6.75, y=1.1, label={[reference]right:\(\bar{b}_6\)}] -- {
                        tbl1s6 / \(\explar{b}{2}{1}\) [x=5, y=0, >shorten >=-3pt],
                        tbl1s7 / \(\explar{b}{3}{2}\) [x=8.5, y=0, >shorten >=-3pt],
                    },
                },
            },
        };
    \end{scope}

    %% Réconciliation T_c
    \begin{scope}[x=.5cm, y=1.1cm, shift={(24, 0)}]
        \graph[no placement, nodes={dependent}]{
            tcl3s1 / [speciation, label={[reference]right:\(\bar{c}_1\)}, x=1.5, y=2] -- {
                tcl1s1 / \(\explar{c}{1}{1}\) [x=0, y=0, >shorten >=-2pt],
                tcl2s1 / [speciation, label={[reference]right:\(\bar{c}_2\)}, x=3, y=1.1] -- {
                    tcl1s2 / \(\explar{c}{2}{1}\) [x=2, y=0, >shorten >=-2pt],
                    tcl1s3 / \(\explar{c}{3}{1}\) [x=4, y=0, >shorten >=-4pt],
                },
            },
        };
    \end{scope}

    \begin{pgfonlayer}{hostnodes}
        \node[host leaf node=color1, fit=(tal1s1)(tal1s4)] (sal1s1) {};
        \node[host leaf node=color1, fit=(tal1s2)(tal1s5)] (sal1s2) {};
        \node[host leaf node=color1, fit=(tal1s3)(tal1s6)] (sal1s3) {};
        \node[host node=color1, fit=(tal2s2)(tal2s3)] (sal2s1) {};
        \node[host node=color1, fit=(tal3s1)(tal3s2)] (sal3s1) {};

        \node[host leaf node=color2, fit=(tbl1s1)(tbl1s2)(tbl1s5), inner xsep=-.15cm] (sbl1s1) {};
        \node[host leaf node=color2, fit=(tbl1s3)(tbl1s6)] (sbl1s2) {};
        \node[host leaf node=color2, fit=(tbl1s4)(tbl1s7)] (sbl1s3) {};
        \node[host node=color2, fit=(tbl2s2)(tbl2s3)] (sbl2s1) {};
        \node[host node=color2, fit=(tbl3s1)(tbl3s2)] (sbl3s1) {};

        \node[host leaf node=color3, fit=(tcl1s1), inner xsep=.05cm] (scl1s1) {};
        \node[host leaf node=color3, fit=(tcl1s2), inner xsep=.05cm] (scl1s2) {};
        \node[host leaf node=color3, fit=(tcl1s3), inner xsep=.05cm] (scl1s3) {};
        \node[host node=color3, fit=(tcl2s1), inner xsep=.15cm] (scl2s1) {};
        \node[host node=color3, fit=(tcl3s1), inner xsep=.15cm] (scl3s1) {};
    \end{pgfonlayer}

    \begin{pgfonlayer}{hostedges}
        \path[host edge=color1] (sal1s1.-10) -- (sal3s1.-10) -- (sal3s1.170) -- (sal1s1.170);
        \path[host edge=color1] (sal1s2.-10) -- (sal2s1.-10) -- (sal2s1.170) -- (sal1s2.170);
        \path[host edge=color1] (sal2s1.25) -- (sal3s1.10) -- (sal3s1.190) -- (sal2s1.190);
        \path[host edge=color1] (sal1s3.10) -- (sal2s1.10) -- (sal2s1.190) -- (sal1s3.190);
        \path[host edge terminal=color1] (sal3s1.east) -- ($(sal3s1.east)+(0, 1)$)
            -- ($(sal3s1.west)+(0, 1)$) -- (sal3s1.west);

        \path[host edge=color2] (sbl1s1.-10) -- (sbl3s1.-10) -- (sbl3s1.170) -- (sbl1s1.170);
        \path[host edge=color2] (sbl1s2.-10) -- (sbl2s1.-10) -- (sbl2s1.170) -- (sbl1s2.170);
        \path[host edge=color2] (sbl2s1.25) -- (sbl3s1.10) -- (sbl3s1.190) -- (sbl2s1.190);
        \path[host edge=color2] (sbl1s3.10) -- (sbl2s1.10) -- (sbl2s1.190) -- (sbl1s3.190);
        \path[host edge terminal=color2] (sbl3s1.east) -- ($(sbl3s1.east)+(0, 1)$)
            -- ($(sbl3s1.west)+(0, 1)$) -- (sbl3s1.west);

        \path[host edge=color3] (scl1s1.-10) -- (scl3s1.-10) -- (scl3s1.170) -- (scl1s1.170);
        \path[host edge=color3] (scl1s2.-10) -- (scl2s1.-10) -- (scl2s1.170) -- (scl1s2.170);
        \path[host edge=color3] (scl2s1.25) -- (scl3s1.10) -- (scl3s1.190) -- (scl2s1.190);
        \path[host edge=color3] (scl1s3.10) -- (scl2s1.10) -- (scl2s1.190) -- (scl1s3.190);
    \end{pgfonlayer}
\end{tikzpicture}

%% Forêt d’adjacences
\vspace{1.5em}
\begin{tikzpicture}[reference/.style={font=\small, inner sep=0pt, outer sep=0pt, outer ysep=3pt}]
    \draw[separator rule] (-3.5, .75) -- ++(14, 0);

    \begin{scope}
        \graph[
            phylogenetic tree layout,
            evolutionary unit length=0.55cm,
            nodes={reference},
        ] {
            / {(\(\bar{a}_1\), \(\bar{b}_1\))} -- {
                / {(\(\bar{a}_3\), \(\bar{b}_3\))} -- {
                    / {(\(\bar{a}_6\), \(\bar{b}_6\))} -- {
                        / {\((\explar{a}{3}{2}, \explar{b}{3}{2})\)},
                        / {\((\explar{a}{2}{2}, \explar{b}{2}{1})\)},
                    },
                    / {\((\explar{a}{1}{2}, \explar{b}{1}{3})\)} [> length=2],
                },
                / {(\(\bar{a}_2\), \(\bar{b}_2\))} -- {
                    / {(\(\bar{a}_5\), \(\bar{b}_5\))} -- {
                        / {\((\explar{a}{3}{1}, \explar{b}{3}{1})\)},
                        / \(\strut\star\),
                    },
                    / {(\(\bar{a}_4\), \(\bar{b}_4\))} -- / {\((\explar{a}{1}{1}, \explar{b}{1}{1})\)},
                },
            },
        };
    \end{scope}

    \begin{scope}[shift={(6, -1.87)}]
        \node[reference] {\((\explar{b}{1}{1}, \explar{b}{1}{2})\)};
    \end{scope}

    \begin{scope}[shift={(9, -.57)}]
        \graph[
            phylogenetic tree layout,
            evolutionary unit length=1.05cm,
            nodes={reference},
        ] {
            / {(\(\bar{b}_6\), \(\bar{c}_2\))} -- {
                / {\((\explar{b}{3}{2}, \explar{c}{3}{1})\)},
                / {\((\explar{b}{2}{1}, \explar{c}{2}{1})\)},
            },
        };
    \end{scope}
\end{tikzpicture}

    \caption[Interpretation of an adjacency reconciliation]{
        \emph{(Top.)}
        Reconciled trees of \(T_a\), \(T_b\) and \(T_c\) with \(S\)
        under the \(\recmodelDL\) model
        and according to the reconciliation represented in \cref{fig:reconciliation:adj-rec}.
        The grey labels make it possible to identify the internal nodes in the lower part of the figure.
        \emph{(Bottom.)}
        Adjacency forest representing a \(\tup{\recmodelDL\rmunit, \recmodelDL\rmadj}\) evolutionary scenario starting from three ancestral adjacencies.
        The star \(\star\) corresponds to a break of the adjacency through the loss of one of its endpoints.
    }
    \label{fig:reconciliation:adj-forest}
\end{figure}

Provided that all events except \(\evtG\) and \(\evtB\) are given a zero cost, \textcite{Berard2012EvolutionGeneNeighborhoods} showed that the \(\tup{\recmodelDL\rmunit, \recmodelDL\rmadj}\)-parsimony problem can be solved in polynomial time by precomputing individual \(\recmodelDL\) reconciliations and then running a dynamic programming algorithm.
No polynomial algorithm is known to date for the general problem and \textcite{Boussau2020ReconcilingGeneTrees} conjectured that the problem is NP-complete.

Following the initial publication of \textcite{Berard2012EvolutionGeneNeighborhoods}, the model was extended by \textcite{Patterson2013LateralGeneTransfer} and then by \textcite{Duchemin2017DecostarReconstructingAncestral} under the names DeCoLT and DeCoSTAR to take into account the \(\evtT\) and \(\evtTX\) events.

Since adjacencies are reconstructed separately from one another, these algorithms can return a non-linear set of adjacencies (see \cref{definition:adjrec-lin}).
A linearization step can be carried out as post-processing, but the linearization problem is NP-complete in general~\parencite{Manuch2012LinearizationAncestralMultichromosomal}.
The non-linearity of an optimal result can, however, help detect inconsistencies in the input data~\parencite{Tannier2020AncestralGenomeOrganization}.

The DeCo approach assumes that adjacencies evolve independently of one another, which excludes from the modelling segmental events affecting several adjacencies simultaneously.
This characteristic also means that the loss of a unit inside a segment of adjacencies causes this segment to split in two, and requires the gain of a new adjacency between the two segments.
A way to remedy this last limitation was, however, recently presented by~\textcite{Tannier2020AncestralGenomeOrganization}.

\subsection{Segmental reconciliation}\label{sec:reconciliation:multiple:segmental}

To allow the modelling of segmental events simultaneously affecting more than two units, \textcite{Delabre2020EvolutionSegmentalDuplications} introduced, under the name \emph{super-reconciliation}, a method generalizing the notion of adjacency to sets of units of arbitrary size called \emph{syntenies}.
The internal order of these syntenies can be taken into account (segments) or ignored (groups).

\begin{definition}[Segments]\label{definition:seg}
    Let \(L\) and \(M\) be two sets and \(f : L \onto M\) a map.
    A \emph{segment} of size~\(k\) is a tuple \(x = (x_1, …, x_k) \in L^k\) such that, for all \(1 \leq i, j \leq k\) with \(i \neq j\), \(x_i \neq x_j\) and \(f(x_i) = f(x_j)\).
    The set of possible segments of arbitrary size is denoted \(\allsegs(f)\).
\end{definition}

\begin{definition}[Groups]\label{definition:cluster}
    Let \(L\) and \(M\) be two sets and \(f : L \onto M\) a map.
    A \emph{group} is a subset \(c = \{c_1, …, c_k\} \subseteq L\) such that, for all \(1 \leq i, j \leq k\), \(f(c_i) = f(c_j)\).
    The set of possible groups is denoted \(\allclstrs(f)\).
\end{definition}

The definitions detailed in the previous section regarding adjacency reconciliations and forests carry over directly to the context of super-reconciliation, to define synteny reconciliations and forests.

This extension makes it possible to specify evolutionary models comprising segmental events, which simultaneously affect sets of units.
The evolutionary model studied in the article builds on \(\recmodelDL\).
During a speciation event~(\(\evtS\)), each member unit of the synteny simultaneously undergoes a codivergence and ends up in the syntenies of both children.
A duplication event~(\(\evtD\)) includes the possibility that only a sub-segment of the initial synteny is copied.
Finally, a loss event~(\(\evtL\)) affects a sub-segment of the initial synteny, forming a new synteny from the concatenation of the two remaining parts.
Formally, for a node \(v \in \tnodes(\segtree)\) of the synteny forest such that \(\exseg(v) = (x_1, …, x_k)\):

\begin{compactitemize}
    \item if \(\evt(\segtree) = \evtS\), then
        \(v\) and all \(x_i\) are binary internal nodes,
        \(\exseg(v\lchild) = (y_1, …, y_k)\) and \(\exseg(v\rchild) = (z_1, …, z_k)\),
        and for all \(1 \leq i \leq k\),
        \(\evt(x_i) = \evtS\)
        and \(\{y_i, z_i\} = \{x_i{}\lchild, x_i{}\rchild\}\);
    \item if \(\evt(\segtree) = \evtD\), then
        \(v\) and all \(x_i\) are also binary,
        and there exist two indices \(1 \leq s \leq t \leq k\) such that
        \(\exseg(v\lchild) = (y_1, …, y_k)\) and \(\exseg(v\rchild) = (z_s, …, z_t)\)
        (or conversely for \(v\lchild\) and \(v\rchild\)),
        and for all \(1 \leq i \leq k\),
        \(\evt(x_i) = \evtD\),
        for \(1 \leq i < s\) and \(t < i \leq k\),
        \(y_i \in \{x_i{}\lchild, x_i{}\rchild\}\),
        and for \(s \leq i \leq t\),
        and \(\{y_i, z_i\} = \{x_i{}\lchild, x_i{}\rchild\}\);
    \item if \(\evt(\segtree) = \evtL\), then
        \(v\) and all \(x_i\) are unary and
        there exist two indices \(1 \leq s \leq t \leq k\) such that
        \(\exseg(v\child) = (y_1, …, y_{s-1}, y_{t+1}, …, y_k)\),
        for all \(1 \leq i < s\) and \(t < i \leq k\),
        \(\evt(x_i) = \evtP\) and
        \(y_i = x_i{}\child\),
        and for \(s \leq i \leq t\),
        \(\evt(x_i) = \evtL\).
\end{compactitemize}

The combinatorial optimization problem associated with super-reconciliation is similar to the \(\tup{\recmodel\rmunit, \recmodel\rmadj}\)-parsimony problem defined previously, except that the expected output is a forest of syntenies instead of adjacencies.
This problem comes in a \enquote{large} version, in which the topology of the forest and the content of its syntenies must be inferred simultaneously, and a \enquote{small} version in which the topology is fixed.

In the article, a restriction of the optimization problem is studied in which the synteny forest is made up of a single tree and all segments derive from a unique ancestral segment.
In its large version, the problem is shown to be NP-complete.
The complexity status of the small version remains unknown in the case of segments, whereas the problem is polynomial in the case of groups.

%! TeX root = ./main.tex
% vim: set spelllang=fr:
\selectlanguage{french}
\section{Comparaison et perspectives}

% TODO: introduire la question de recherche ici

Dans ce chapitre ont été présentées deux approches généralisant la réconciliation bilatérale à des contextes où plusieurs unités dépendantes évoluent conjointement.
Ce cas de figure se présente lorsque de multiples gènes évoluent ensemble le long d’un génome.
Le principe général consiste à intégrer un niveau intermédiaire dans la hiérarchie évolutive, décrivant l’évolution d’adjacences (DeCo) ou de segments porteurs d’ensembles d’unités (super-réconciliation).

L’approche DeCo permet de modéliser des événements qui affectent des paires d’unités voisines, incluant notamment des réarrangements.
D’un point de vue algorithmique, elle consiste à réconcilier individuellement les arbres phylogénétiques de chaque unité, puis à rassembler des paires de nœuds pour former une forêt d’adjacences.
Chaque racine de cette forêt correspond au gain d’une adjacence qui se propage ensuite dans sa descendance.
L’objectif est de minimiser le nombre de racines, c’est-à-dire le nombre de gains d’adjacences.
Le statut de complexité du problème est encore inconnu, quoiqu’il ait été conjecturé NP-complet, mais un algorithme polynomial existe lorsque les réconciliations individuelles sont fixées en avance.
La modélisation à l’échelle des adjacences impose des limites sur les événements pouvant être pris en compte, puisqu’elle exclut les événements qui affectent simultanément plus de deux unités.
Par ailleurs, un résultat optimal peut comporter des structures non-linéaires, dans lesquelles une unité est voisine de trois autres ou plus, incompatibles avec les structures linéaires ou circulaires des génomes.

La super-réconciliation propose un paradigme similaire, basé sur l’évolution de segments joignant un nombre quelconque d’unités.
Cette extension permet de modéliser des événements segmentaux affectant plusieurs unités.
De la même façon que pour DeCo, l’objectif algorithmique est de construire un arbre dans lequel chaque nœud rassemble des segments de nœuds issus de réconciliations individuelles.
La structure reconstruite ici est un arbre et non une forêt car tous les segments sont supposés dériver d’un segment ancestral unique.
Le problème d’optimisation combinatoire associé est NP-complet dans sa version générale, lorsque la topologie de l’arbre de synténies est inconnue.
Si la topologie est fixée, la complexité est inconnue dans la déclinaison du problème avec des segments, et polynomiale dans sa déclinaison avec des groupes.
Les modèles d’évolution considérés jusqu’à présent pour la super-réconciliation sont encore limités.
Contrairement à ceux considérés sous DeCo, ils ne permettent par exemple pas de modéliser des duplications en tandem, des fusions ou fissions de synténies ou des réarrangements.

\subsection{Thèse}

L’objectif général du projet de recherche est de développer la théorie de la réconciliation multiple, en particulier autour de l’approche de super-réconciliation, et de démontrer son potentiel dans l’analyse de l’évolution de systèmes biologiques.

% Contributions

